% ============================================================
\documentclass[master.tex]{subfiles}

\begin{document}

% ------------------------------------------------------------
%  Title block (standalone) / chapter heading (master)
% ------------------------------------------------------------
\ifSubfilesClassLoaded{%
  \begin{center}
    {\LARGE\bfseries\color{isublue}
     Chapter 7: Instrumental Variables}\\[6pt]
    {\large Theory and Methods in Causal Inference}\\[4pt]
    {\normalsize Department of Statistics, Iowa State University}
  \end{center}
  \bigskip
  \hrule height 1.2pt \relax
  \smallskip
}{%
  \chapter{Instrumental Variables}
}

% ------------------------------------------------------------
%  Learning objectives
% ------------------------------------------------------------
\begin{tcolorbox}[defstyle, title={Learning Objectives}]
By the end of this chapter, students should be able to:
\begin{enumerate}
  \item Diagnose the endogeneity problem: explain why an unobserved
    confounder makes back-door adjustment fail, and describe how an
    instrumental variable provides an alternative identification
    route.
  \item State the three IV assumptions---relevance, exogeneity, and
    exclusion---in the graphical, potential-outcomes, and
    econometric languages, and explain which are testable and which
    require institutional justification.
  \item Follow how the IV assumptions identify the Wald estimand,
    and interpret the reduced form and first stage as its two
    observable components.
  \item Explain what the Wald estimand identifies when treatment
    effects are heterogeneous---the Local Average Treatment Effect
    for compliers---why that estimand depends on the choice of
    instrument, and when it coincides with the ATE.
  \item Compare IV and back-door adjustment on the assumptions each
    requires, the estimand each identifies, and the ways each can
    fail.
\end{enumerate}
\end{tcolorbox}


% ============================================================
\section{Why Instrumental Variables?}
\label{w7:sec:why}
% ============================================================

Chapter~6 showed how causal effects can be identified when all
confounders are observed and can be blocked by conditioning on~$X$.
When some confounders are unobserved, back-door adjustment fails.
This chapter develops instrumental variables (IV) as an alternative
identification strategy: rather than blocking the confounding path
$T \leftarrow U \rightarrow Y$, IV exploits an external variable~$Z$
whose effect on~$T$ is free of confounding by~$U$.  This chapter
asks what IV identifies and under what assumptions; Chapter~12 asks
how that estimand is computed and tested in practice.

\subsection{The Endogeneity Problem}

Chapter~6 established that the propensity score provides a powerful
surrogate for randomization, but only under \emph{unconfoundedness}:
\[
  \bigl(Y(0),\, Y(1)\bigr) \indep T \;\big|\; X.
\]
This assumption requires that every variable affecting both treatment
and outcome is observed and included in~$X$.  In many empirical
settings this is implausible:
\begin{itemize}
  \item In labor economics, unobserved ability or motivation affects
    both schooling decisions and wages.
  \item In epidemiology, unobserved health behaviors affect both
    treatment uptake and outcomes.
  \item In survey sampling, voluntary participation is driven by
    factors that also affect the variables of interest.
\end{itemize}
Whenever an unobserved confounder~$U$ creates a back-door path
$T \leftarrow U \rightarrow Y$, the adjustment formula fails:
\[
  \int f(y \mid t, x)\, p(x)\, dx
  \;\neq\;
  f\!\bigl(y \mid \doop(T{=}t)\bigr).
\]
The gap is the \emph{endogeneity bias}.  We need a different
identification strategy.

\subsection{The IV Idea}

An \emph{instrumental variable} $Z$ is an observed variable that:
\begin{enumerate}
  \item \emph{Moves} treatment~$T$ (relevance).
  \item \emph{Does so exogenously}---$Z$ is unrelated to the
    unobserved confounder~$U$ (exogeneity).
  \item \emph{Affects~$Y$ only through~$T$}---$Z$ has no direct
    path to~$Y$ (exclusion).
\end{enumerate}
Under these three conditions, the variation in~$T$ induced by~$Z$
is free of confounding by~$U$.  The ratio of the $Z$-induced
variation in~$Y$ to the $Z$-induced variation in~$T$ identifies
the causal effect.

Instrumental variables do not eliminate the confounding path
$T \leftarrow U \rightarrow Y$, nor do they make treatment as-if
randomly assigned for the full population.  Instead, IV identifies
causal effects by isolating a component of treatment variation that
is induced by~$Z$ and is therefore exogenous under the IV
assumptions.  In that sense, IV \emph{avoids} confounding rather
than controlling for it---a distinction that matters both for
interpreting what is identified (the LATE for compliers, not the
ATE) and for understanding which assumption does the heaviest
lifting (exclusion, not unconfoundedness).

\begin{example}{Returns to Schooling as Motivation}
\label{w7:ex:schooling-motivation}
A researcher wants to estimate the causal effect of years of
schooling~$T$ on log wages~$Y$.  Unobserved ability~$U$ raises
both schooling and wages, so $\mathrm{Cov}(T, \varepsilon) > 0$
and OLS is biased upward.  No set of observed covariates~$X$
fully captures ability.  An instrumental variable~$Z$ that shifts
schooling for reasons unrelated to ability---such as proximity to
a school or a policy change in compulsory attendance laws---provides
a way to isolate the causal effect of schooling on wages.
\end{example}

\begin{figure}[h]
\centering
\begin{tikzpicture}[
  node distance = 12mm and 18mm,
  w7obs/.style  = {circle, draw=accent, line width=1pt,
                   minimum size=9mm, font=\small, fill=defbg},
  w7unobs/.style= {circle, draw=backdoorred, line width=1pt, dashed,
                   minimum size=9mm, font=\small, fill=warnbg},
  w7xobs/.style = {circle, draw=causalgreen, line width=1pt,
                   minimum size=9mm, font=\small,
                   fill=green!8},
  w7arr/.style  = {-{Stealth[length=4.5pt]}, line width=0.9pt,
                   color=accent},
  w7redarr/.style={-{Stealth[length=4.5pt]}, line width=0.9pt,
                   color=backdoorred, dashed},
  w7lbl/.style  = {font=\footnotesize\color{isublue}}
]
%% ── Left panel: Back-door approach ─────────────────────────────────
\begin{scope}[xshift=0cm]
  \node[w7obs]                          (T1) {$T$};
  \node[w7obs, right=of T1]             (Y1) {$Y$};
  \node[w7unobs, above=of T1, xshift=9mm] (U1) {$U$};
  \node[w7xobs, below=of T1, xshift=9mm]  (X1) {$X$};

  \draw[w7arr]    (T1) -- (Y1);
  \draw[w7redarr] (U1) -- (T1);
  \draw[w7redarr] (U1) -- (Y1);
  \draw[-{Stealth[length=4.5pt]}, line width=0.9pt, color=causalgreen]
    (X1) -- (T1);
  \draw[-{Stealth[length=4.5pt]}, line width=0.9pt, color=causalgreen]
    (X1) -- (Y1);

  \node[w7lbl, below=18mm of T1, xshift=9mm, align=center]
    {\textbf{Back-door approach}\\[2pt]
     Condition on $X$ (observed)\\to block
     $T \leftarrow U \rightarrow Y$};
\end{scope}
%% ── Right panel: IV approach ────────────────────────────────────────
\begin{scope}[xshift=7.2cm]
  \node[w7obs]                   (Z2) {$Z$};
  \node[w7obs, right=of Z2]      (T2) {$T$};
  \node[w7obs, right=of T2]      (Y2) {$Y$};
  \node[w7unobs, above=of T2, xshift=9mm] (U2) {$U$};

  \draw[w7arr]    (Z2) -- node[above, w7lbl]{\scriptsize IV} (T2);
  \draw[w7arr]    (T2) -- (Y2);
  \draw[w7redarr] (U2) -- (T2);
  \draw[w7redarr] (U2) -- (Y2);

  \node[w7lbl, below=18mm of T2, align=center]
    {\textbf{IV approach}\\[2pt]
     Use only $Z$-induced variation in $T$\\
     ($Z \indep U$, so this variation is clean)};
\end{scope}
\end{tikzpicture}

\medskip
{\small Dashed red arrows denote paths through the unobserved
confounder~$U$.  In the back-door approach, conditioning on the
observed~$X$ (green) blocks the confounding path.  In the IV
approach, $Z$ is used to select only the part of $T$'s variation
that is free of~$U$.}
\label{w7:fig:bd-vs-iv}
\end{figure}


% ============================================================
\section{Graphical Setup and Core Assumptions}
\label{w7:sec:assumptions}
% ============================================================

This section is the conceptual core of the chapter.  We fix the
causal graph for the IV model, translate the three identifying
assumptions into each of the three languages used in this course,
and explain why the do-calculus formulation is the most precise.
Mastering these assumptions in all three languages is essential:
the potential outcomes language dominates biostatistics and program
evaluation, the structural language is standard in econometrics,
and the do-calculus language is the course's preferred vehicle for
identification proofs.

\subsection{The IV DAG}

The causal structure for a basic IV model with covariates~$X$ is:

\begin{center}
\begin{tikzpicture}[
  node distance = 14mm and 20mm,
  w7obs/.style = {circle, draw=accent, line width=1pt,
                  minimum size=10mm, align=center, font=\small,
                  fill=defbg},
  w7unobs/.style={circle, draw=backdoorred, line width=1pt, dashed,
                  minimum size=10mm, align=center, font=\small,
                  fill=warnbg},
  w7arr/.style  = {-{Stealth[length=5pt]}, line width=1pt, color=accent},
  w7redarr/.style={-{Stealth[length=5pt]}, line width=1pt,
                   color=backdoorred, dashed},
  w7lbl/.style  = {font=\footnotesize\itshape\color{isublue}}
]
  \node[w7obs]               (Z) {$Z$};
  \node[w7obs, right=of Z]   (T) {$T$};
  \node[w7obs, right=of T]   (Y) {$Y$};
  \node[w7obs, below=12mm of T] (X) {$X$};
  \node[w7unobs, above right=10mm and 10mm of T] (U) {$U$};

  \draw[w7arr]    (Z) -- node[above, w7lbl]{(1)} (T);
  \draw[w7arr]    (T) -- (Y);
  \draw[w7arr]    (X) -- (T);
  \draw[w7arr]    (X) -- (Y);
  \draw[w7arr]    (X) -- (Z);
  \draw[w7redarr] (U) -- (T);
  \draw[w7redarr] (U) -- (Y);
\end{tikzpicture}
\end{center}

\medskip\noindent The three IV conditions correspond to three
features of this DAG:
\begin{enumerate}
  \item \textbf{Relevance}: there is a directed path $Z \to T$
    (edge labeled~(1) above).
  \item \textbf{Exogeneity}: there is no back-door path from $Z$
    to~$Y$ via~$U$---the DAG implies $Z \indep U \mid X$ by
    $d$-separation.
  \item \textbf{Exclusion}: there is no direct arrow from $Z$
    to~$Y$---$Z$ does not appear in the structural equation
    for~$Y$.
\end{enumerate}

\noindent\textbf{What the DAG rules out.}  The absence of a direct
$Z \to Y$ arrow encodes the exclusion restriction.  Any such
arrow---no matter how small the coefficient---would violate exclusion
and invalidate the Wald estimand.  The DAG forces the researcher to
commit to this assumption explicitly; the potential outcomes
framework leaves it implicit in the notation
$Y_i(t,z) = Y_i(t,z')$.

\subsection{The Three Assumptions in Three Languages}

The same three assumptions can be expressed in three causal
languages.  These are parallel formulations, not literally identical
statements: each language encodes a different aspect of the
underlying causal claim.  The assumptions are ordered
\emph{relevance $\to$ exogeneity $\to$ exclusion}, reflecting the
natural sequence in which a researcher assesses them: first check
that the instrument actually moves treatment, then verify it is
exogenously assigned, and finally---most subtly---verify it has no
direct effect on the outcome.

\medskip
\renewcommand{\arraystretch}{2.0}
\begin{tabular}{@{}p{2.1cm} p{3.9cm} p{3.9cm} p{3.9cm}@{}}
  \toprule
  \textbf{Assumption}
    & \textbf{Potential outcomes}
    & \textbf{Structural / econometric}
    & \textbf{Do-calculus / graphical} \\
  \midrule
  \textbf{Relevance}
    & $Z \nindep T$
    & $\mathrm{Cov}(Z, T) \neq 0$
    & No $d$-separation of $Z$ and $T$ in $\Gcal$ \\
  \textbf{Exogeneity}
    & $Z \indep (Y(0), Y(1), T(0), T(1)) \mid X$
    & $\E[\varepsilon \mid Z, X] = 0$
    & $Z \indep U \mid X$ \\
  \textbf{Exclusion}
    & $Y_i(t, z) = Y_i(t, z')$ for all $z, z'$
    & $Z$ absent from structural equation for $Y$
    & $f(y \mid x, \doop(T{=}t), z)
       = f(y \mid x, \doop(T{=}t))$ \\
  \bottomrule
\end{tabular}
\renewcommand{\arraystretch}{1}

\begin{remark}[Three Languages, One Identification Argument]
These three formulations are aligned but not literally identical.
The graphical statement encodes causal structure---which arrows are
present or absent in~$\Gcal$.  The potential-outcomes statement
encodes counterfactual independence---which potential outcomes are
independent of~$Z$.  The econometric statement encodes moment
orthogonality---which inner product with the residual is zero.
In a well-specified IV model they support the same identification
argument, but students should not treat them as interchangeable
symbols.

Substantively: \emph{relevance} is about the $Z \to T$ link---$Z$
must genuinely move~$T$, not merely be correlated with it.
\emph{Exogeneity} is about the absence of omitted common causes
linking~$Z$ to~$Y$---$Z$ must be as-if randomly assigned relative
to the latent determinants of the outcome.  \emph{Exclusion} is
about the absence of any direct causal path $Z \to Y$ that
bypasses~$T$---it rules out both a direct arrow in the DAG and any
parallel channel through unmodeled variables.
\end{remark}

\medskip\noindent\textbf{Why the do-calculus column is preferred.}
The exclusion restriction in the do-calculus column reads:
\[
  f\!\bigl(y \mid x,\, \doop(T{=}t),\, z\bigr)
  \;=\;
  f\!\bigl(y \mid x,\, \doop(T{=}t)\bigr).
\]
This is a statement about the \emph{interventional} density---the
distribution of~$Y$ after we have \emph{set} $T = t$ by do-surgery.
It says that once~$T$ is fixed by intervention, $Z$ carries no
additional information about~$Y$.  The do-operator makes it
impossible to confuse the exclusion restriction with the
observational statement
$f(y \mid x, T{=}t, z) = f(y \mid x, T{=}t)$,
which is a much weaker condition.

\subsection{Relevance}

\begin{defn}{Relevance}
\label{w7:defn:relevance}
The instrument~$Z$ is \textbf{relevant} if it has a non-zero causal
effect on the treatment~$T$, possibly after conditioning on
covariates~$X$:
\[
  \mathrm{Cov}(Z, T \mid X) \neq 0.
\]
Equivalently, $Z$ and~$T$ are not $d$-separated in the DAG~$\Gcal$.
Relevance is a requirement on the $Z \to T$ \emph{causal} link, not
merely on the observational correlation; the latter can be non-zero
even when the causal effect is zero, due to a common cause of~$Z$
and~$T$.
\end{defn}

Relevance is the population condition that the instrument changes
treatment status: $P(T \mid Z{=}z)$ must vary with~$z$, or in
linear models, $\pi \neq 0$.  The first-stage $F$-statistic is a
sample diagnostic for whether this link is strong enough to support
reliable estimation; it is evidence about instrument strength, not
the assumption itself.  What happens when relevance holds only
weakly---or fails entirely---is analyzed in
\S\ref{w7:sec:why-matter} and \S\ref{w7:sec:overid}.

\subsection{Exogeneity}

\begin{defn}{Exogeneity (Graphical Assumption)}
\label{w7:defn:exogeneity}
The instrument~$Z$ is \textbf{exogenous} given covariates~$X$ if
there is no open back-door path from~$Z$ to~$Y$ through unobserved
causes.  In DAG terms: $Z \indep U \mid X$ in~$\Gcal$ by
$d$-separation.
\end{defn}

\noindent Exogeneity means that the instrument carries no information
about latent determinants of the outcome other than through treatment.
In the DAG language, there is no open back-door path from~$Z$ to~$Y$
through unobserved causes~$U$: the graphical condition $Z \indep U
\mid X$ in~$\Gcal$ captures this directly.  In the structural linear
model, the same requirement appears as the moment orthogonality
condition $\E[Z\varepsilon \mid X] = 0$, because the structural
residual~$\varepsilon$ is a function of~$U$.  In potential-outcomes
language, $Z$ must be independent of the counterfactual quantities
that the IV argument requires---outcomes $Y(0), Y(1)$ and treatment
indicators $T(0), T(1)$---given~$X$.  These are related
manifestations of the same identification requirement, but they live
in different formal languages and are not literally the same
statement.

\noindent\textbf{Econometric implication.}  The moment condition
$\E[\varepsilon Z \mid X] = 0$ is a \emph{consequence} of the
graphical assumption, not a definition of exogeneity.  This
distinction matters: a researcher who adopts the moment condition as
a primitive has no guarantee that~$Z$ is free of back-door paths to
the outcome, only that a particular linear projection is zero.
The condition enters the Wald derivation in \S\ref{w7:sec:linear}
specifically because it follows from the graphical assumption.

\noindent\textbf{Potential-outcomes analogue.}  In the
potential-outcomes language, exogeneity is expressed as the full
independence of the instrument from all counterfactual outcomes and
counterfactual treatments:
\[
  Z \indep \bigl(Y(0),\, Y(1),\, T(0),\, T(1)\bigr) \;\big|\; X.
\]
This is a stronger statement than the moment condition alone: it
also covers the joint independence of~$Z$ with the treatment
potential outcomes $T(0), T(1)$, which the LATE proof in
\S\ref{w7:sec:late} requires in addition to
$\E[\varepsilon Z \mid X] = 0$.

\begin{example}{Quarter of Birth as an Instrument}
\label{w7:ex:qob}
\citet{angrist1991does} used quarter of birth as an instrument
for schooling to estimate the return to education.  In graphical
terms, exogeneity requires no open back-door path from quarter of
birth to wages through unobserved variables such as ability or
family background; in potential-outcomes terms, it requires that
birth timing is independent of what wages a student would earn
under any schooling level.  Both are plausible because birth timing
is largely outside parental control.  \citet{bound1995problems}
later raised concerns about weak first stages in some
specifications, illustrating that even a well-motivated instrument
can fail the relevance requirement in certain samples.
\end{example}

\noindent\textbf{Key point.}  All three versions of exogeneity
are untestable: the unobserved confounder~$U$ is, by definition,
unobserved, and counterfactual potential outcomes are never jointly
observed.  The assumption must be justified by institutional
knowledge or the design of the study.
See~\S\ref{w7:sec:overid} for the limited consistency check that
overidentification provides.

\subsection{Exclusion}

\begin{defn}{Exclusion Restriction}
\label{w7:defn:exclusion}
The instrument~$Z$ satisfies the \textbf{exclusion restriction} if
it affects the outcome~$Y$ \emph{only} through its effect on the
treatment~$T$.  In do-calculus terms:
\[
  f\!\bigl(y \mid x,\, \doop(T{=}t),\, z\bigr)
  \;=\;
  f\!\bigl(y \mid x,\, \doop(T{=}t)\bigr)
  \quad \text{for all } z.
\]
In potential outcomes terms: $Y_i(t, z) = Y_i(t, z')$ for all
$z \neq z'$.
\end{defn}

\noindent Exclusion is a \emph{causal} restriction, not a
conditional-correlation restriction.  It is emphatically
\emph{not} the observational statement $Y \indep Z \mid T, X$,
which can hold in the data even when~$Z$ has a direct structural
effect on~$Y$.  What exclusion requires is that changing~$Z$
while holding treatment fixed---that is, intervening to
set~$T = t$---would leave~$Y$ unchanged.  In the DAG this means
no direct arrow $Z \to Y$; in the structural linear model it means
no~$Z$ term appears in the outcome equation for~$Y$.

\begin{warn}{Why Exclusion Is the Hardest Assumption}
Relevance can be tested.  Exogeneity can sometimes be defended by
design (e.g., randomized encouragement).  But the exclusion
restriction---that $Z$ has \emph{no direct effect} on $Y$
whatsoever---is:
\begin{itemize}
  \item \textbf{Untestable} in just-identified models (one
    instrument, one treatment).
  \item \textbf{Easy to violate in practice.}  An instrument that
    ``moves treatment'' often also moves other inputs.  For example:
    if~$Z$ is distance to a hospital (instrument for treatment
    uptake), it may also directly affect health outcomes through
    travel time in emergencies.
  \item \textbf{Consequential.}  Even a small direct effect of~$Z$
    on~$Y$---if correlated with the instrument---can cause
    substantial bias in the Wald estimand, especially when the
    first stage is weak.  This bias is derived explicitly in
    \S\ref{w7:sec:why-matter}.
\end{itemize}
\end{warn}

\noindent\textbf{The do-calculus makes the violation visible.}
Adding a direct arrow $Z \to Y$ to the DAG is exactly the formal
counterpart of exclusion failing: the arrow represents a structural
effect of~$Z$ on~$Y$ that bypasses~$T$.  In the mutilated graph
$\Gcal_{\overline{T}}$ (where we have removed all arrows into~$T$),
the path $Z \to Y$ remains open.  Rule~1 of the do-calculus, which would allow us to remove~$Z$
from the density of~$Y$ given $\doop(T{=}t)$, no longer applies.
The exclusion restriction is therefore visibly a statement about the
structure of $\Gcal_{\overline{T}}$, not of~$\Gcal$ itself.


\begin{remark}[Three Assumptions Are Necessary but Not Sufficient for Point Identification]
\label{w7:rem:three-not-enough}
The three IV assumptions---relevance, exogeneity, and exclusion---are
necessary but not sufficient for \emph{point} identification of any
causal estimand.  Under these three assumptions alone, the observed
data restrict the distribution of potential outcomes but do not pin it
down to a single value; the Wald ratio is a well-defined function of
observable quantities, but it does not equal any specific causal
parameter without additional structure
\citep{hernan2006instruments, levis2024nonparametric}.
A fourth \emph{structural assumption} on the counterfactual
distribution of treatment response or treatment effect heterogeneity
is required.

The choice of fourth assumption determines which causal quantity the
Wald estimand identifies.  Two leading choices appear in this chapter:
constant treatment effects (Framework~1, \S\ref{w7:sec:linear}), under
which the Wald estimand identifies the ATE; and monotonicity
(Framework~2, \S\ref{w7:sec:late}), under which it identifies the LATE
for compliers.  A third important family relies on \emph{structural
mean model} (SMM) restrictions.  Under an additive SMM that rules
out direct effect modification of~$Y$ by~$Z$ among the treated,
the Wald ratio identifies the ATT---the average effect of treatment
among those who actually received it---without requiring either
constant effects or monotonicity \citep{hernan2006instruments,
levis2024nonparametric}.  Moving further to identify the population
ATE from the ATT requires an additional step: that the average effect
among the treated equals the average effect among the untreated, a
condition that follows from constant effects but must otherwise be
separately justified.

Strikingly, across these different identification schemes the same
conditional Wald formula
$\mathrm{Cov}(Z, Y \mid X)/\mathrm{Cov}(Z, T \mid X)$ appears as
the identifying expression in many cases.  The qualification
``many'' is precise: when identification relies on a
\emph{multiplicative} SMM instead of an additive one, the identifying
formula for the ATE takes a different form entirely and the Wald ratio
is no longer consistent for the target estimand
\citep{hernan2006instruments}.  What changes in the additive cases is
not the observed-data formula but its causal target: the same
estimator can consistently estimate the ATE, the LATE, or the ATT
depending on which structural assumption is maintained.  This is
one of the clearest illustrations in causal inference of why the
choice of causal model---not just the choice of estimator---determines
what is being learned from data.

A further consequence deserves emphasis: the three core IV assumptions
are all expressible within the graphical language of this course.
Relevance is the directed edge $Z \to T$; exogeneity is the absence
of an open back-door path from~$Z$ to~$Y$ via~$U$; exclusion is the
absence of a direct edge $Z \to Y$ in the mutilated graph
$\Gcal_{\overline{T}}$.  The fourth assumption, by contrast, lies
\emph{outside} the graphical language in every case.  Monotonicity
($T_i(1) \geq T_i(0)$ for all~$i$) is a restriction on the joint
distribution of potential treatments across two intervention values
simultaneously---a cross-world statement that cannot be encoded by
the presence or absence of any arrow in the DAG.  Constant treatment
effects ($Y_i(1) - Y_i(0) = \beta$) similarly constrains the
covariance structure of the potential outcome pair $(Y_i(0), Y_i(1))$,
which the DAG does not represent.  This is a genuine scope limitation
of the graphical approach: a DAG specifies the causal structure of a
single world, but the fourth assumption always reaches across worlds.
The practical implication is that IV identification cannot be
completed by causal graph analysis alone; the fourth assumption must
be supplied from substantive knowledge about the mechanism, and
defended on those grounds.
\end{remark}

\begin{example}{Same Observables, Different ATEs}
\label{w7:ex:ident-failure}
The following construction makes Remark~\ref{w7:rem:three-not-enough}
concrete.  Two distinct data-generating processes satisfy all three IV
assumptions and produce identical observable statistics, yet imply
different ATEs.  No amount of data can distinguish them without a
fourth structural assumption.

\medskip
\noindent\textbf{Observable statistics.}
Let $Z, T \in \{0,1\}$ and let $Y$ be a continuous outcome.
Suppose the joint distribution of $(Y, T, Z)$ satisfies:
\begin{align*}
  \E[T \mid Z{=}1] - \E[T \mid Z{=}0] &= 0.2 \quad
    \text{(first stage)}, \\
  \E[Y \mid Z{=}1] - \E[Y \mid Z{=}0] &= 0.1 \quad
    \text{(reduced form)},
\end{align*}
giving a Wald ratio of $0.1/0.2 = 0.5$.  Both DGPs share the same
compliance structure---under monotonicity (no defiers),
$P(\mathrm{co}) = 0.2$, $P(\mathrm{at}) = 0.4$,
$P(\mathrm{nt}) = 0.4$---though under DGP~A the compliance types
are simply a convenient labeling; what drives identification there is
the constant-effect restriction, not monotonicity.

\medskip
\noindent\textbf{DGP~A: Constant treatment effects.}
Every unit has the same effect $\tau_i = 0.5$.  Then:
\[
  \text{reduced form}
  = P(\mathrm{co}) \cdot \tau_{\mathrm{co}}
  = 0.2 \times 0.5 = 0.1 \;\checkmark
  \qquad
  \mathrm{ATE} = 0.5.
\]

\medskip
\noindent\textbf{DGP~B: Heterogeneous treatment effects with
monotonicity.}  Effects differ by compliance type:
$\tau_{\mathrm{co}} = 0.5$, $\tau_{\mathrm{at}} = 0$,
$\tau_{\mathrm{nt}} = 0$.  Then:
\[
  \text{reduced form}
  = P(\mathrm{co}) \cdot \tau_{\mathrm{co}}
  = 0.2 \times 0.5 = 0.1 \;\checkmark
  \qquad
  \mathrm{ATE}
  = 0.2 \times 0.5 + 0.4 \times 0 + 0.4 \times 0
  = 0.1.
\]

\medskip
\noindent\textbf{What this shows.}
Both DGPs produce the same first stage, the same reduced form, and
therefore the same Wald ratio of $0.5$.  Yet the ATE is $0.5$ under
DGP~A and $0.1$ under DGP~B---a fivefold difference.  The three IV
assumptions do not distinguish between them.  The \emph{fourth}
assumption is what resolves the ambiguity: under DGP~A, the constant
effects assumption makes the Wald ratio equal to the ATE; under
DGP~B, the monotonicity assumption makes the Wald ratio equal to the
LATE ($= 0.5$) for compliers, which is a different quantity from the
ATE ($= 0.1$).  Committing to one fourth assumption or the other
changes not the estimator but what it estimates.
\end{example}

% ============================================================
\section{Identification in the Linear Homogeneous-Effect Model}
\label{w7:sec:linear}
% ============================================================

\begin{tcolorbox}[defstyle, title={Framework 1: Constant Treatment Effect and Linear Structure}]
This section works entirely within a linear structural model in
which the treatment effect is the \emph{same} for every unit:
$Y_i(t) - Y_i(0) = \beta$ for all~$i$.  Under this assumption, IV
identifies the single parameter~$\beta$, which equals the ATE, the
ATT, and the LATE simultaneously.  Section~\ref{w7:sec:late}
relaxes this restriction and introduces Framework~2.

\smallskip
\noindent\textbf{Note on assumption strength.}  Constant treatment
effects is the strongest form of the homogeneity family of fourth
assumptions: it requires not only that unobserved confounders~$U$ do
not modify the effect of~$T$ on~$Y$, but that the effect is exactly
the same for every unit within each stratum of~$X$.  Weaker
homogeneity conditions---such as requiring only that~$U$ not act as
an additive effect modifier on~$Y$---can also identify the ATE under
IV \citep{wang2018bounded}.
\end{tcolorbox}

\subsection{The Linear Structural Model}

We begin with the homogeneous-effect linear SEM because it makes the
IV logic transparent.  In Section~\ref{w7:sec:late} we relax
homogeneity and reinterpret the same Wald ratio as a LATE, but the
identification argument is easier to see first in this simpler
setting.

Consider the linear structural model:
\begin{align}
  Y &= \alpha + \beta T + \gamma^\top X + \varepsilon,
  \label{w7:eq:outcome} \\
  T &= \pi Z + \delta^\top X + \eta,
  \label{w7:eq:firststage}
\end{align}
where $\varepsilon$ and $\eta$ are structural errors with
$\mathrm{Cov}(\varepsilon, \eta) = \rho\sigma\tau \neq 0$.  The
non-zero covariance is the source of endogeneity: OLS applied
to~\eqref{w7:eq:outcome} gives a biased estimator of~$\beta$.

The parameter of interest is~$\beta$, the causal effect of~$T$
on~$Y$:
\[
  \beta = \E\!\bigl[Y \mid \doop(T{=}t+1)\bigr]
        - \E\!\bigl[Y \mid \doop(T{=}t)\bigr].
\]
Equation~\eqref{w7:eq:outcome} is the \emph{outcome equation} (or
second stage); equation~\eqref{w7:eq:firststage} is the
\emph{first-stage} equation.  The \emph{reduced form} is obtained by
substituting~\eqref{w7:eq:firststage} into~\eqref{w7:eq:outcome}:
\[
  Y = \alpha + \beta(\pi Z + \delta^\top X + \eta)
    + \gamma^\top X + \varepsilon
  = \alpha + \beta\pi Z + (\beta\delta + \gamma)^\top X
    + (\beta\eta + \varepsilon).
\]
The reduced-form coefficient on~$Z$ is $\beta\pi$: the total
effect of the instrument on the outcome.  Dividing by the
first-stage coefficient~$\pi$ recovers~$\beta$, provided
$\pi \neq 0$.

\subsection{The OLS Bias}

To see why OLS fails, apply OLS to~\eqref{w7:eq:outcome}.  The
probability limit of the OLS estimator is:
\[
  \mathrm{plim}\; \hat\beta_{\mathrm{OLS}}
  = \beta + \frac{\mathrm{Cov}(T, \varepsilon)}{\mathrm{Var}(T)}
  = \beta + \frac{\rho\sigma\tau}{\pi^2 \mathrm{Var}(Z) + \tau^2}.
\]
The second term is the endogeneity bias.  It is zero only if
$\rho = 0$ (no unobserved confounding) or if the instrument
perfectly determines~$T$ (infeasible in practice).  The direction
of the bias depends on the sign of~$\rho$.

The point of IV is not that OLS is always wrong, but that when~$T$
is endogenous, IV can recover a causal parameter from the exogenous
variation in~$T$ that~$Z$ induces---variation that OLS mixes
indiscriminately with the confounded component.

\subsection{Derivation of the Wald Estimand}

\begin{tcolorbox}[warnstyle, title={Assumptions Active in This Derivation}]
\begin{enumerate}
  \item \textbf{Linear homogeneous-effect model:} $Y_i(t) - Y_i(0)
    = \beta$ for all~$i$, and the structural equations are linear.
  \item \textbf{Relevance:} $\mathrm{Cov}(Z, T \mid X) \neq 0$.
  \item \textbf{Exogeneity:} $\mathrm{Cov}(Z, \varepsilon \mid X)
    = 0$.
  \item \textbf{Exclusion:} $Z$ has no direct effect on~$Y$; it
    does not appear in the structural equation for~$Y$.
\end{enumerate}
Assumptions 2--4 are the three IV conditions; Assumption~1 is the
homogeneity restriction specific to this section.  Under
heterogeneous effects the same Wald ratio identifies a different
quantity; see Section~\ref{w7:sec:late}.
\end{tcolorbox}

We derive the identification of~$\beta$ step by step from the three
IV assumptions, using the structural model as the vehicle and the
do-calculus to interpret the exclusion step.
\begin{enumerate}
  \item \textbf{Exogeneity} ($Z \indep U \mid X$ in $\Gcal$) implies,
    for the structural residual~$\varepsilon$:
    \[
      \E[\varepsilon \mid Z, X] = \E[\varepsilon \mid X].
    \]
    With the normalization $\E[\varepsilon \mid X] = 0$ (absorbed
    into~$\alpha$), we obtain $\E[\varepsilon \mid Z, X] = 0$.

  \item \textbf{Exclusion} ($Z$ absent from the structural equation
    for~$Y$) means that~$\varepsilon$ depends on~$U$ but not
    on~$Z$ directly.  Combined with exogeneity, this gives:
    \[
      \E[\varepsilon \cdot Z \mid X] = 0.
    \]
    In do-calculus terms: in $\Gcal_{\overline{T}}$, Rule~1 removes
    $Z$ from $f(y \mid x, \doop(T{=}t), z)$, because
    $Y \indep Z \mid T, X$ in $\Gcal_{\overline{T}}$.

  \item Taking expectations of~\eqref{w7:eq:outcome} after
    multiplying through by $(Z - \E[Z \mid X])$ and applying
    step~2:
    \begin{align*}
      \E\!\bigl[(Y - \alpha - \gamma^\top X)
               (Z - \E[Z\mid X])\bigr]
      &= \beta\, \E\!\bigl[(T - \E[T\mid X])
                            (Z - \E[Z\mid X])\bigr]
         + \underbrace{\E[\varepsilon (Z - \E[Z\mid X])]}_{=\,0}.
    \end{align*}
    In the simpler case with no covariates~$X$, this collapses to
    the \textbf{reduced-form decomposition}:
    \[
      \mathrm{Cov}(Y, Z) = \beta\, \mathrm{Cov}(T, Z),
    \]
    which shows that the $Z$--$Y$ covariance is entirely
    attributable to the causal path $Z \to T \to Y$ (by
    exclusion), and equals~$\beta$ times the first-stage
    covariance.

  \item \textbf{Relevance} ($\mathrm{Cov}(Z, T \mid X) \neq 0$)
    ensures the denominator is non-zero, so we can solve:
    \[
      \beta
      = \frac{\mathrm{Cov}(Y,\, Z \mid X)}
             {\mathrm{Cov}(T,\, Z \mid X)}.
    \]
\end{enumerate}

In the binary instrument case ($Z \in \{0,1\}$) without additional
covariates~$X$, this simplifies to the \textbf{Wald estimand}:

\begin{tcolorbox}[defstyle, title={The Wald Estimand}]
\label{w7:eq:wald}
\[
  \beta
  \;=\;
  \frac{\E[Y \mid Z{=}1] - \E[Y \mid Z{=}0]}
       {\E[T \mid Z{=}1] - \E[T \mid Z{=}0]}.
\]
The numerator is the \emph{reduced form}: the total effect of~$Z$
on~$Y$.  The denominator is the \emph{first stage}: the effect
of~$Z$ on~$T$.  The ratio removes the portion of the reduced form
that passes through channels other than~$T$---but only because
the exclusion restriction guarantees there are none.
\end{tcolorbox}

At this point the Wald estimand identifies~$\beta$ only because the
model assumes a common treatment effect across all units.  Under
heterogeneous effects, the same ratio generally has a different
interpretation; see Section~\ref{w7:sec:late}.

\begin{remark}[Estimation of the Wald Estimand]
The Wald estimand above is an identification result: it expresses the
causal parameter~$\beta$ as a ratio of observable quantities.
Estimation---that is, how to consistently estimate this ratio from
a finite sample, including the correct treatment of standard errors
---is taken up in Chapter~12.
\end{remark}

\begin{example}{Returns to Schooling}
\label{w7:ex:returns}
In the Mincer earnings equation
$\log W = \alpha + \beta S + \gamma^\top X + \varepsilon$,
where~$S$ is years of schooling, $W$ is wages, and~$X$ are
demographics, the OLS estimator of~$\beta$ is biased upward because
unobserved ability~$U$ raises both~$S$ and~$W$ ($\rho > 0$).
\citet{angrist1991does} use quarter of birth as~$Z$: proximity
to mandatory school-leaving age at different birth quarters generates
exogenous variation in completed schooling.  The IV estimate of the
return to schooling is approximately $0.08$--$0.10$ per year,
similar to the OLS estimate in this case, suggesting the ability
bias is modest.
\end{example}


% ============================================================
\section{Why the IV Assumptions Matter}
\label{w7:sec:why-matter}
% ============================================================

The three IV assumptions are not merely technical conditions: each
one is load-bearing, and each failure mode produces a distinct,
quantifiable distortion of the Wald estimand.  Violating exogeneity
or exclusion contaminates the reduced form, and dividing by a weak
first stage magnifies that contamination.  A weak instrument is
therefore not merely inefficient; it makes any small violation of
the identifying assumptions more consequential.

\medskip\noindent\textbf{When relevance fails.}
If $\pi = 0$, the Wald estimand is undefined: its denominator is
zero and there is no exogenous variation to exploit.  When~$\pi$ is
small but nonzero, the instrument is \emph{weak}.  The estimator's
variance diverges as $\pi \to 0$ because the denominator amplifies
noise, and finite-sample bias pulls the IV estimate toward the OLS
estimate at a rate proportional to~$1/F$, where~$F$ is the
first-stage $F$-statistic.

\medskip\noindent\textbf{When exclusion fails.}
If the outcome equation contains a direct effect
$Y = \alpha + \beta T + \delta Z + \gamma^\top X + \varepsilon$
with $\delta \neq 0$, the reduced form picks up both the
indirect path~$\beta\pi$ and the direct path~$\delta$, so the
Wald estimand converges to $\beta + \delta/\pi$.  The
bias~$\delta/\pi$ is amplified by a weak first stage: a small
direct effect combined with a weak instrument can produce large
bias.  This is why a weak instrument with a plausible exclusion
violation is not ``nearly valid''---it may be severely misleading.

\medskip\noindent\textbf{When exogeneity fails.}
If $\mathrm{Cov}(Z, \varepsilon \mid X) \neq 0$, the IV moment
condition fails and the Wald estimand converges to
$\beta + \mathrm{Cov}(\varepsilon, Z)/\mathrm{Cov}(T, Z)$.
This is the IV analogue of OLS omitted-variable bias: IV fails
when the instrument itself is endogenous, just as OLS fails when
the treatment is endogenous.  Again the bias is amplified when
the first stage is weak.

\medskip
Table~\ref{w7:tab:testability} summarizes the testability status
of the three assumptions.  The key asymmetry is that the hardest
work in defending an IV design---the institutional arguments for
exogeneity and exclusion---is precisely the work that the data
cannot do for the researcher.

\begin{table}[h]
\centering
\renewcommand{\arraystretch}{1.6}
\begin{tabular}{p{2.4cm} p{3.4cm} p{6.0cm}}
  \toprule
  \textbf{Assumption} & \textbf{Directly testable?} & \textbf{Basis for assessment} \\
  \midrule
  Relevance
    & Yes
    & First-stage $F$-statistic (diagnostic for instrument strength,
      not the assumption itself) \\
  Exogeneity
    & No
    & Institutional knowledge; randomization
      (if available); placebo regressions on
      pre-determined outcomes \\
  Exclusion
    & No in just-identified case;
      partially in overidentified case
    & Institutional argument; overidentification
      test ($J$-test, Chapter~12) checks mutual
      consistency of instruments but cannot confirm all
      are valid \\
  \bottomrule
\end{tabular}
\renewcommand{\arraystretch}{1}
\caption{Testability of the three IV assumptions.  Although exogeneity
  and exclusion are not directly testable---no finite dataset can
  confirm them---they are not entirely beyond empirical scrutiny.
  The IV model imposes inequality restrictions on the observable
  data distribution; when these are violated, falsification tests can
  detect the violation \citep{kitagawa2015test, hernan2006instruments}.
  Importantly, however, such tests are not consistent against all
  alternatives: for most data distributions under which the assumptions
  fail, the tests will not reject even in large samples.  The
  distinction between ``directly testable,'' ``falsifiable in some
  cases,'' and ``requires institutional argument'' captures the
  practical hierarchy for IV assumption assessment.}
\label{w7:tab:testability}
\end{table}




% ============================================================
\section{Lab: OLS vs.\ IV Across Instrument Strengths}
\label{w7:sec:lab}
% ============================================================

The analytic results of Sections~\ref{w7:sec:linear}
and~\ref{w7:sec:why-matter} established two facts: OLS is
biased by a precise, computable amount whenever
$\rho \neq 0$; and IV is consistent but its variance diverges
as $\pi \to 0$.  What the analytics do not immediately reveal
is the \emph{finite-sample consequence}: for weak instruments,
the IV variance penalty is so severe that the biased OLS
estimator can have lower mean squared error.  This lab
verifies the bias formulas numerically and traces the
bias--variance tradeoff across the full range of instrument
strength, from $\pi = 0$ (no instrument) to $\pi = 2$
(very strong instrument).  This simulation studies estimator
behavior conditional on the IV assumptions being true; it does
not address the separate and fundamentally substantive question
of whether a proposed instrument is valid in an applied study.

% ----------------------------------------
\subsection*{Data-Generating Process}
\addcontentsline{toc}{subsection}{Data-Generating Process}

\begin{tcolorbox}[defstyle, title={DGP for Lab~7}]
\label{w7:lab:dgp}
The simulation implements the linear structural model of
Section~\ref{w7:sec:linear} with no covariates~$X$.  Each
replication draws $n = 500$ observations from:
\begin{align}
  Y_i &= \beta T_i + \varepsilon_i, \label{w7:lab:outcome} \\
  T_i &= \pi Z_i + \eta_i,          \label{w7:lab:fs}
\end{align}
with $Z_i \overset{\mathrm{i.i.d.}}{\sim} \mathrm{N}(0,1)$ and
the structural errors generated as:
\[
  \eta_i = U_i, \qquad
  \varepsilon_i = \rho U_i + \sqrt{1 - \rho^2}\,\xi_i,
  \qquad U_i,\, \xi_i \overset{\mathrm{i.i.d.}}{\sim} \mathrm{N}(0,1).
\]
Here $U_i$ is the unobserved confounder shared by both equations,
and $\xi_i$ is the idiosyncratic component of~$\varepsilon_i$
that is independent of everything else.
This gives $\mathrm{Var}(\varepsilon_i) = \mathrm{Var}(\eta_i) = 1$
and $\mathrm{Cov}(\varepsilon_i, \eta_i) = \rho$, matching
the notation $\rho\sigma\tau$ from Section~\ref{w7:sec:linear}
with $\sigma = \tau = 1$.  The fixed parameter values are:
\[
  \beta = 1 \text{ (true causal effect)},
  \qquad
  \rho = 0.8 \text{ (strong positive endogeneity)}.
\]
The first-stage coefficient~$\pi$ is varied across eight
values to sweep from a completely uninformative instrument
to a very strong one:
\[
  \pi \;\in\; \{0,\; 0.10,\; 0.15,\; 0.20,\; 0.30,\; 0.50,\; 1.00,\; 2.00\}.
\]
The expected first-stage $F$-statistic is approximately
$n\pi^2 = 500\pi^2$ (since $Z \sim \mathrm{N}(0,1)$ and
$\mathrm{Var}(\eta) = 1$).  The instrument is undefined at
$\pi = 0$; the Staiger--Stock rule of thumb $F \geq 10$
corresponds to $\pi \geq 0.141$ in this design.
\end{tcolorbox}

\begin{tcolorbox}[defstyle, title={Potential Outcomes Interpretation of the DGP}]
\label{w7:lab:po}
The structural DGP above has an exact translation into the
potential outcomes language of Chapter~5, which makes explicit
\emph{why} OLS fails and \emph{why} IV succeeds.

\medskip\noindent\textbf{Potential outcomes.}
The latent background factor $\varepsilon_i$ is fixed for each
unit across all counterfactual worlds.  The potential outcome
under treatment value~$t$ is therefore:
\begin{equation}
  Y_i(t) \;=\; \beta t + \varepsilon_i.
  \label{w7:lab:po-eq}
\end{equation}
The observed outcome satisfies $Y_i = Y_i(T_i) = \beta T_i +
\varepsilon_i$, which is the consistency assumption.  Because
$\varepsilon_i$ does not depend on~$t$, the individual treatment
effect is constant:
\[
  Y_i(t) - Y_i(t') \;=\; \beta(t - t') \quad \text{for all } i.
\]
This is exactly Framework~1 from Section~\ref{w7:sec:linear}:
$\beta$ is simultaneously the ATE, the ATT, and the LATE.

\medskip\noindent\textbf{The unobserved confounder.}
The shared factor $U_i$ is the unobserved confounder (plotted
as the dashed node in the DAG of
Section~\ref{w7:sec:assumptions}).  It enters both the
treatment assignment and the background outcome:
\[
  T_i = \pi Z_i + U_i,
  \qquad
  \varepsilon_i = \rho U_i + \sqrt{1-\rho^2}\,\xi_i.
\]
Because $\varepsilon_i$ depends on $U_i$ and $\eta_i = U_i$,
units with large~$T_i$ (driven by large~$U_i$) also tend to
have large~$\varepsilon_i$, and hence large~$Y_i$, \emph{even
if $\beta = 0$}.  Formally,
\[
  \mathrm{Cov}\!\bigl(Y_i(t),\, T_i\bigr)
  = \mathrm{Cov}(\varepsilon_i,\, \eta_i)
  = \mathrm{Cov}(\rho U_i,\, U_i) = \rho \;\neq\; 0,
\]
so the unconfoundedness condition $(Y_i(t) \indep T_i)$
fails.  The parameter $\rho$ is precisely the
\emph{degree of violation}: when $\rho = 0$ there is no
confounding and OLS is unbiased; when $\rho = 0.8$ the
confounder $U_i$ explains $80\%$ of the co-movement
between~$T$ and~$Y$ that is not due to $\beta$.

\medskip\noindent\textbf{Why IV restores identification.}
The instrument $Z_i$ is drawn independently of $U_i$
(and of $\xi_i$), so it is independent of both the
unobserved confounder and the background outcome:
\[
  Z_i \;\perp\!\!\!\perp\; U_i
  \quad\Longrightarrow\quad
  Z_i \;\perp\!\!\!\perp\; \varepsilon_i.
\]
This is exogeneity, satisfied by construction.  The
instrument $Z_i$ does not appear in $Y_i(t)$
(equation~\eqref{w7:lab:po-eq}), which is the exclusion
restriction: $Y_i(t, z) = Y_i(t)$ for all~$z$.  And
$Z_i$ moves~$T_i$ through the first-stage coefficient~$\pi$,
which is relevance.  All three IV assumptions hold exactly
by construction, confirming that the Wald estimator targets
the true $\beta$ for any $\pi \neq 0$.

\medskip\noindent\textbf{What varies across the grid.}
The potential outcomes $Y_i(t)$ are \emph{the same} for
every value of~$\pi$ --- the parameter $\pi$ governs only
how strongly the instrument shifts treatment.  Varying $\pi$
changes the \emph{first-stage signal} but not the
counterfactual surfaces.  All that changes for OLS as
$\pi$ grows is that more of $T_i$'s variance comes from the
clean source~$Z_i$ rather than from the confounded
source~$U_i$, which partially dilutes the endogeneity
bias.  But because the potential outcomes themselves are
unchanged, $\beta = 1$ remains the target at every row of
Table~\ref{w7:tab:lab}.
\end{tcolorbox}

% ----------------------------------------
\subsection*{Estimators}
\addcontentsline{toc}{subsection}{Estimators}

Two estimators of $\beta$ are compared at each value of $\pi$.

\medskip\noindent\textbf{OLS.}  Regress $Y$ on $T$:
\[
  \hat\beta_{\mathrm{OLS}}
  = \frac{\mathrm{Cov}(Y, T)}{\mathrm{Var}(T)}.
\]
From Section~\ref{w7:sec:linear}, this converges in probability to
\[
  \mathrm{plim}\;\hat\beta_{\mathrm{OLS}}
  = \beta + \frac{\rho}{\pi^2 + 1}
  = 1 + \frac{0.8}{\pi^2 + 1}.
\]
The bias is always positive (since $\rho > 0$), decreasing
in $|\pi|$, but never zero for finite $\pi$.

\medskip\noindent\textbf{IV (Wald estimator).}  Use $Z$ as
the instrument:
\[
  \hat\beta_{\mathrm{IV}}
  = \frac{\mathrm{Cov}(Y, Z)}{\mathrm{Cov}(T, Z)}.
\]
This estimator is consistent for $\beta$ for any $\pi \neq 0$.
Its finite-sample variance is approximately
$\sigma^2 / (n \pi^2)$, diverging as $\pi \to 0$.

% ----------------------------------------
\subsection*{Results}
\addcontentsline{toc}{subsection}{Results}

Table~\ref{w7:tab:lab} reports Monte Carlo mean, bias, standard
deviation, and RMSE across $B = 2{,}000$ replications at
$n = 500$, seed \texttt{2024}, alongside the theoretical OLS
bias from Section~\ref{w7:sec:linear} and the mean first-stage
$F$-statistic.

\begin{table}[h]
\centering
\renewcommand{\arraystretch}{1.5}
\caption{Monte Carlo performance of OLS and IV across instrument
  strengths.  True $\beta = 1.000$.  Theoretical OLS bias $=
  0.8/(\pi^2 + 1)$.  $n = 500$, $B = 2{,}000$ replications,
  seed \texttt{2024}.  ``---'' denotes $\pi = 0$: IV is not
  defined.  IV mean and SD at $\pi = 0.10$ and $0.15$ are
  dominated by extreme outliers; the IV median is $\approx 1.01$
  at both values.}
\label{w7:tab:lab}
\setlength{\tabcolsep}{5pt}
\begin{tabular}{crr rrrr rrrr}
\toprule
& & & \multicolumn{4}{c}{\textbf{OLS}} &
      \multicolumn{4}{c}{\textbf{IV (Wald)}} \\
\cmidrule(lr){4-7}\cmidrule(lr){8-11}
$\pi$ & $F$ & \shortstack{Theory\\bias} &
  Mean & Bias & SD & RMSE &
  Mean & Bias & SD & RMSE \\
\midrule
$0.00$ &    $1$ & $+0.800$ & $1.801$ & $+0.801$ & $0.026$ & $0.802$ &
  \multicolumn{4}{c}{---} \\
$0.10$ &    $6$ & $+0.792$ & $1.791$ & $+0.791$ & $0.027$ & $0.792$ &
  $0.872$ & $-0.128$ & $4.391$ & $4.392$ \\
$0.15$ &   $13$ & $+0.782$ & $1.782$ & $+0.782$ & $0.027$ & $0.783$ &
  $0.764$ & $-0.236$ & $6.487$ & $6.491$ \\
$0.20$ &   $21$ & $+0.769$ & $1.769$ & $+0.769$ & $0.027$ & $0.770$ &
  $0.940$ & $-0.060$ & $0.380$ & $0.385$ \\
$0.30$ &   $46$ & $+0.734$ & $1.735$ & $+0.735$ & $0.028$ & $0.735$ &
  $0.975$ & $-0.025$ & $0.163$ & $0.165$ \\
$0.50$ &  $127$ & $+0.640$ & $1.640$ & $+0.640$ & $0.028$ & $0.640$ &
  $0.996$ & $-0.004$ & $0.091$ & $0.091$ \\
$1.00$ &  $500$ & $+0.400$ & $1.401$ & $+0.401$ & $0.027$ & $0.402$ &
  $0.999$ & $-0.001$ & $0.045$ & $0.045$ \\
$2.00$ & $2006$ & $+0.160$ & $1.160$ & $+0.160$ & $0.019$ & $0.161$ &
  $1.000$ & $\phantom{-}0.000$ & $0.022$ & $0.022$ \\
\bottomrule
\end{tabular}
\end{table}

\medskip
Four lessons stand out.

\medskip
\noindent\textbf{1.\ The OLS bias formula is exact.}
The ``Theory bias'' column ($0.8/(\pi^2 + 1)$) matches the
simulated OLS bias to four decimal places across all eight
values of~$\pi$.  OLS is biased in the direction of~$\rho$ (positive
here) at every value of~$\pi$, including~$\pi = 0$ where there
is no instrument at all.  Crucially, the OLS bias \emph{decreases}
as~$\pi$ grows, because a stronger first stage means more of
$T$'s variance comes from the exogenous source~$Z$ rather than
from the endogenous error~$\eta$.  At~$\pi = 2$ the bias is only
$0.160$, yet this seemingly small bias still produces RMSE $= 0.161$,
which is $7\times$ larger than the IV RMSE of $0.022$ at the
same~$\pi$.

\medskip
\noindent\textbf{2.\ IV is consistent but has catastrophically heavy
tails when the instrument is weak.}
At $\pi = 0.10$ (F $\approx 6$) and $\pi = 0.15$ (F $\approx 13$)
the IV \emph{mean} is far from the true value ($0.87$ and $0.76$
respectively), suggesting bias.  This is misleading: the IV
\emph{median} is $\approx 1.01$ at both values, confirming that
the estimator is consistent and the \emph{typical} replication
recovers~$\beta = 1$.  The mean is dragged off by a small
fraction of replications in which the first stage happens to be
near zero, making the Wald ratio explode.  These tail events
inflate the standard deviation to $4.4$ and $6.5$
respectively---orders of magnitude larger than the OLS SD of
$0.027$.  The take-away is that the weak instrument problem is
not inconsistency; it is uncontrolled tail behavior, and the
mean and RMSE are highly sensitive to it.

\medskip
\noindent\textbf{3.\ The RMSE crossover occurs near $F \approx 20$.}
Comparing the RMSE columns: OLS RMSE ranges from $0.802$
($\pi = 0$) down to $0.161$ ($\pi = 2$), while IV RMSE starts
infinite ($\pi = 0$), is catastrophic at $\pi = 0.10$--$0.15$,
then drops sharply to $0.385$ at $\pi = 0.20$ (F $\approx 21$)
and to $0.165$ at $\pi = 0.30$ (F $\approx 46$).  IV first beats
OLS on RMSE at $\pi = 0.20$ ($0.385 < 0.770$): this is the
crossover.  Below this threshold, the variance penalty of IV
exceeds the bias penalty of OLS and the biased estimator is
\emph{preferable} in mean squared error terms.
The Staiger--Stock rule of thumb ($F \geq 10$, i.e.\
$\pi \geq 0.141$ in this design) is therefore slightly
\emph{too lenient}: at F $\approx 13$ IV RMSE is still $6.5$,
nearly $8\times$ OLS.  A more conservative threshold of
$F \geq 20$--$25$ is needed before IV dominates OLS in this DGP.

\medskip
\noindent\textbf{4.\ A strong instrument eliminates both problems.}
At $\pi = 1.00$ (F $\approx 500$), IV is virtually unbiased
(bias $= -0.001$) with SD $= 0.045$ and RMSE $= 0.045$,
while OLS still carries a bias of $0.401$, giving RMSE $= 0.402$
--- a \textbf{9-fold} improvement from IV.  At $\pi = 2.00$
(F $\approx 2006$) the improvement is $7\times$ ($0.022$ vs
$0.161$).  With a strong instrument, IV is simultaneously unbiased
\emph{and} more efficient than OLS in MSE terms, because OLS
efficiency is illusory: its small variance is offset by a large,
persistent bias.

\begin{warn}{RMSE Cannot Be the Only Criterion}
Table~\ref{w7:tab:lab} shows that OLS has lower RMSE than IV
for $\pi < 0.20$ in this DGP.  This does not vindicate OLS.
An estimator with RMSE $= 0.78$ because it is biased by $0.80$
is useless for causal inference: the bias is systematic and
does not shrink with sample size.  As $n \to \infty$, IV RMSE
shrinks to zero while OLS RMSE stays at $0.80$.  The RMSE
comparison is only meaningful in \emph{finite samples}.  For
a fixed~$n$, it quantifies when a weak instrument is so
unreliable that IV is not yet practically useful --- but that
is an argument for finding a stronger instrument, not for
using OLS.
\end{warn}

\begin{remark}[The First-Stage $F$-Statistic as a Diagnostic]
\label{w7:rem:fstat}
The approximate formula $F \approx n\pi^2$ (which the simulation
confirms to within Monte Carlo error throughout) gives
first-stage $F$-statistics of $6$, $13$, $21$, $46$, $127$,
$500$, and $2006$ for the seven non-zero $\pi$ values in
Table~\ref{w7:tab:lab}.  The Staiger--Stock threshold of $F \geq 10$
\citep{bound1995problems} is widely used in practice and
corresponds to IV RMSE within a factor of roughly two of the
strong-instrument limit.  The simulation here suggests this
threshold may understate the problem when endogeneity is strong
($\rho = 0.8$): at F $= 13$ the IV RMSE is $6.5$, far
above any useful threshold.  A researcher should report the
first-stage $F$ alongside any IV estimate, and treat values
below $20$--$25$ with particular caution.
\end{remark}

% ============================================================
\section{Multiple Instruments and Overidentification}
\label{w7:sec:overid}
% ============================================================

This section gives only the identification-level intuition for using
more than one instrument.  The computational details of 2SLS with
multiple instruments, GMM weighting, and the Sargan--Hansen $J$-test
are deferred to Chapter~12.

When there are exactly as many instruments as endogenous variables
($q = p$), the model is \emph{just-identified}: the IV assumptions
pin down the causal parameter exactly, leaving no surplus
identifying variation.  When $q > p$, the model is
\emph{overidentified}: the extra instruments impose additional
moment restrictions beyond those needed for identification.  Under
exogeneity and exclusion, every valid instrument must imply the
same causal estimate, so those extra restrictions are testable.
The Sargan--Hansen $J$-test \citep{sargan1958estimation,hansen1982large}
checks whether the instruments are mutually consistent in this
sense; rejection indicates that at least one instrument violates
exogeneity or exclusion, but does not identify which one.  The
test statistic and its asymptotic distribution are derived in
Chapter~12.

\begin{warn}{What Overidentification Does Not Do}
Passing the $J$-test does not confirm that any instrument is valid.
It only shows that the sample moment conditions are mutually
compatible---a weaker conclusion than validity.  In the
just-identified case the test has no power at all.  Multiple
invalid instruments can agree with one another if they share the
same violation, so a consistent set of instruments is not
necessarily a valid one.  Heterogeneous treatment effects or
model misspecification can also affect the test's behavior.
Overidentification is an opportunity for a consistency check, not
a substitute for the institutional argument that makes a design
credible.
\end{warn}

\begin{remark}[GMM and Overidentification]
In the overidentified case, the surplus moment conditions can also
be exploited for efficient estimation: rather than discarding the
extra instruments, one can combine all $q$ moment conditions
optimally.  The generalized method of moments (GMM) framework
provides the natural setting for this, treating the just-identified
IV estimator as a special case.  This connection is developed in
Chapter~12.
\end{remark}


% ============================================================
\section{Heterogeneous Treatment Effects and the LATE Framework}
\label{w7:sec:late}
% ============================================================

\begin{tcolorbox}[defstyle, title={Framework 2: Heterogeneous Treatment Effects}]
Section~\ref{w7:sec:linear} identified the causal effect under the
assumption that every unit shares the same treatment effect~$\beta$.
In reality, treatment effects are typically heterogeneous:
$\tau_i = Y_i(1) - Y_i(0)$ varies across units.  This section
drops the homogeneity assumption entirely and asks: what does the
Wald estimand identify when effects are heterogeneous?  The answer
is the \emph{Local Average Treatment Effect} (LATE) --- the average
effect for the specific subpopulation whose treatment status changes
with the instrument.
\end{tcolorbox}

\subsection{Compliance Types}

In the binary instrument, binary treatment setting
($Z, T \in \{0,1\}$), treatment effects may be
\emph{heterogeneous}: different units may respond differently to
treatment.  To understand \emph{what} IV identifies, we classify
units by how their treatment decision responds to the instrument.

\begin{defn}{Compliance Types \citep{angrist1996identification}}
\label{w7:defn:compliance}
For binary $Z, T \in \{0,1\}$, define the potential treatment
$T_i(z)$ as the treatment unit~$i$ would take under instrument
value~$z$.  The four compliance types are:

\medskip
\begin{center}
\begin{tabular}{lcc}
  \toprule
  \textbf{Compliance type} & $T_i(0)$ & $T_i(1)$ \\
  \midrule
  Complier     & 0 & 1 \\
  Always-taker & 1 & 1 \\
  Never-taker  & 0 & 0 \\
  Defier       & 1 & 0 \\
  \bottomrule
\end{tabular}
\end{center}
\medskip

\noindent A \textbf{complier} takes treatment if and only if the
instrument is switched on.  An \textbf{always-taker} takes
treatment regardless of~$Z$.  A \textbf{never-taker} never takes
treatment regardless of~$Z$.  A \textbf{defier} does the opposite
of what the instrument suggests.
\end{defn}

The instrument~$Z$ only shifts treatment for \emph{compliers}:
always-takers and never-takers have the same treatment status
regardless of~$Z$, so they contribute nothing to the denominator
$\E[T \mid Z{=}1] - \E[T \mid Z{=}0]$.  Intuitively, the Wald
estimand is driven entirely by the complier subgroup.
Compliance type is defined by a pair of potential treatment statuses
$(T_i(0), T_i(1))$, so it is a latent causal classification rather
than an observed covariate-defined subgroup.

\subsection{The Monotonicity Assumption}

Under the first three IV assumptions alone, the Wald estimand does
not have a clean causal interpretation when treatment effects are
heterogeneous, because defiers can contaminate the estimand.  An
additional assumption eliminates them.

\begin{defn}{Monotonicity \citep{imbens1994identification}}
\label{w7:defn:monotonicity}
The treatment assignment is \textbf{monotone in}~$Z$ if:
\[
  T_i(1) \;\geq\; T_i(0) \quad \text{for all } i.
\]
Equivalently: there are no defiers in the population.
\end{defn}

\noindent\textbf{Interpretation.}  Monotonicity says that switching
the instrument from~$0$ to~$1$ never \emph{discourages} treatment.
If~$Z$ is a randomized encouragement, it is plausible that some
units respond (compliers) and others do not (always-takers or
never-takers), but that no unit actively avoids treatment
\emph{because of} the encouragement.  Monotonicity is most
credible when the instrument is a designed encouragement or
allocation rule.  In many observational IV settings, however, the
no-defiers assumption is substantively much harder to defend and
requires explicit justification.

\subsection{The LATE Theorem}

We now drop the homogeneous-effect assumption of
Section~\ref{w7:sec:linear}.  The same Wald ratio remains
identified, but it generally no longer equals the ATE; instead it
identifies the average treatment effect for compliers.

\begin{thm}{LATE Theorem \citep{imbens1994identification}}
\label{w7:thm:late}
Under the four IV assumptions (relevance, exogeneity, exclusion, and
monotonicity), the Wald estimand identifies the \textbf{Local
Average Treatment Effect} (LATE):
\[
  \frac{\E[Y \mid Z{=}1] - \E[Y \mid Z{=}0]}
       {\E[T \mid Z{=}1] - \E[T \mid Z{=}0]}
  \;=\;
  \E\!\bigl[Y(1) - Y(0) \;\big|\; T_i(1) > T_i(0)\bigr]
  \;\equiv\; \tau_{\mathrm{LATE}}.
\]
The Wald estimand identifies the average causal effect for
\emph{compliers} only.
\end{thm}

\begin{proof}
We decompose the numerator and denominator by compliance type.

\medskip\noindent\textbf{Denominator.}
\begin{align*}
  \E[T \mid Z{=}1] - \E[T \mid Z{=}0]
  &= \E[T(1)] - \E[T(0)]
   = \E[T(1) - T(0)]
   = P(\text{complier}),
\end{align*}
using exogeneity ($Z \indep T(0), T(1)$) in the first equality,
and monotonicity in the last: always-takers contribute $1 - 1 = 0$,
never-takers contribute $0 - 0 = 0$, and defiers do not exist,
so only compliers (with $T(1) - T(0) = 1$) contribute:
\[
  \E[T \mid Z{=}1] - \E[T \mid Z{=}0] = P(\text{complier}).
\]

\medskip\noindent\textbf{Numerator.}  By exogeneity and the law
of total expectation, conditioning on compliance type:
\begin{align*}
  \E[Y \mid Z{=}1] - \E[Y \mid Z{=}0]
  &= \E[Y(T(1))] - \E[Y(T(0))] \\
  &= \sum_{c} P(c)\,
     \E\!\bigl[Y(T_c(1)) - Y(T_c(0))\bigr].
\end{align*}
For always-takers: $T(1) = T(0) = 1$, contribution $= 0$.
For never-takers: $T(1) = T(0) = 0$, contribution $= 0$.
For compliers: $T(1) = 1, T(0) = 0$, so
$Y(T(1)) - Y(T(0)) = Y(1) - Y(0)$.
No defiers exist by monotonicity.  Therefore:
\[
  \E[Y \mid Z{=}1] - \E[Y \mid Z{=}0]
  = P(\text{complier})\,
    \E\!\bigl[Y(1) - Y(0) \;\big|\; \text{complier}\bigr].
\]

\medskip\noindent\textbf{Ratio.}  Dividing numerator by denominator:
\[
  \frac{\E[Y \mid Z{=}1] - \E[Y \mid Z{=}0]}
       {\E[T \mid Z{=}1] - \E[T \mid Z{=}0]}
  = \E\!\bigl[Y(1) - Y(0) \;\big|\; \text{complier}\bigr]
  = \tau_{\mathrm{LATE}}. 
\]
\end{proof}

\begin{remark}[The Two Senses of ``Local'']
The LATE is local in two senses: it is local to the complier group
defined by this instrument, and local to the particular instrument
that defines that group.  A different instrument, even for the same
treatment, will generally select a different complier population and
identify a different LATE.  This is the subject of
Section~\ref{w7:subsec:diff-instr}.
\end{remark}


% ============================================================
\section{Interpreting IV Estimands}
\label{w7:sec:interpret}
% ============================================================

The two frameworks---linear homogeneous effects
(\S\ref{w7:sec:linear}) and the LATE framework
(\S\ref{w7:sec:late})---give the Wald estimand different
interpretations.  This section synthesizes those interpretations
and clarifies when they agree, when they disagree, and what the
difference means for applied work.

\subsection{What the Two Frameworks Say}

\begin{tcolorbox}[defstyle, title={What IV Identifies: A Comparison of the Two Frameworks}]
\renewcommand{\arraystretch}{1.6}
\begin{tabular}{p{2.4cm} p{4.6cm} p{5.0cm}}
  \toprule
  & \textbf{Framework~1}
    (linear, homogeneous)
  & \textbf{Framework~2}
    (heterogeneous effects) \\
  \midrule
  Key assumption
    & $\tau_i = \beta$ for all~$i$
    & Monotonicity; no defiers \\
  What IV identifies
    & $\beta = \mathrm{ATE} = \mathrm{ATT} = \mathrm{LATE}$
    & $\tau_{\mathrm{LATE}} = \E[\tau_i \mid \text{complier}]$ \\
  Estimand depends on instrument?
    & No (same~$\beta$ regardless of~$Z$)
    & Yes (different~$Z$ $\Rightarrow$ different compliers
      $\Rightarrow$ different LATE) \\
  Required for ATE?
    & Yes, automatically
    & Only if all units are compliers or effects homogeneous \\
  \bottomrule
\end{tabular}
\renewcommand{\arraystretch}{1}
\end{tcolorbox}

Framework~1 is a special case of Framework~2: when
$\tau_i = \beta$ for all~$i$, the LATE equals the ATE equals~$\beta$,
and the instrument does not affect the estimand---only identification.
Framework~2 is the more general and realistic setting.  In applied
work, the default interpretation of the Wald estimand is the LATE;
the ATE interpretation requires the additional homogeneity argument
of Framework~1.

\begin{remark}[Same Formula, Different Estimand]
\label{w7:rem:same-formula}
The two frameworks illustrate a broader principle in IV analysis.
Across a range of fourth assumptions---constant effects, monotonicity,
structural mean model restrictions, and others---the conditional Wald
formula $\mathrm{Cov}(Z, Y \mid X)/\mathrm{Cov}(Z, T \mid X)$ serves
as the identifying expression in many (though not all) cases
\citep{levis2024nonparametric}.  Accordingly, two researchers who use
the same instrument and compute the same Wald ratio may be consistently
estimating different causal quantities if they maintain different
structural assumptions: one estimates the ATE (under constant effects),
another estimates the LATE (under monotonicity), and a third estimates
the ATT (under a structural mean model restriction).  The point
estimates and confidence intervals can be numerically identical; what
differs is the causal quantity they represent.

This is not a deficiency of IV but a precise illustration of why
causal modeling is not merely a formal preliminary to estimation.  The
data alone cannot resolve which estimand the Wald ratio identifies;
that determination requires the researcher to commit to a structural
assumption about treatment response.  Stating that commitment
explicitly---and defending it on substantive grounds---is as important
as any statistical calculation.
\end{remark}

\subsection{When Does LATE Equal ATE?}

LATE equals ATE only under additional structure, most notably
treatment-effect homogeneity or special forms of heterogeneity
that make complier effects representative of the full population.
Neither condition should be assumed without argument.

In general, $\tau_{\mathrm{LATE}} \neq \mathrm{ATE}$ when treatment
effects are heterogeneous and some units are not compliers.  To see
this, decompose the ATE by compliance type:
\[
  \mathrm{ATE}
  = P(\mathrm{co})\,\E[\tau_i \mid \mathrm{co}]
  + P(\mathrm{at})\,\E[\tau_i \mid \mathrm{at}]
  + P(\mathrm{nt})\,\E[\tau_i \mid \mathrm{nt}],
\]
where $\mathrm{co}$, $\mathrm{at}$, and $\mathrm{nt}$ abbreviate
complier, always-taker, and never-taker.  The LATE equals only the
first term divided by its probability weight.  The ATE and LATE
coincide if and only if:
\begin{enumerate}
  \item Treatment effects are homogeneous
    ($\E[\tau_i \mid \mathrm{co}] = \E[\tau_i \mid \mathrm{at}]
     = \E[\tau_i \mid \mathrm{nt}]$), or
  \item Everyone is a complier ($P(\mathrm{at}) = P(\mathrm{nt})
     = 0$), which would require~$Z$ to perfectly determine~$T$,
     or
  \item The average effects for always-takers and never-takers
    happen to equal the LATE---an untestable coincidence.
\end{enumerate}
In practice, none of these conditions is likely to hold exactly.
The LATE is a well-defined, identifiable parameter, but it is not
the ATE.

\subsection{Different Instruments, Different Estimands}
\label{w7:subsec:diff-instr}

Because the LATE is specific to the complier population, and
different instruments select different complier populations, two
valid instruments for the same treatment can legitimately identify
different LATEs.  This is a precise, substantive fact---not a
contradiction.  Different instruments can legitimately identify
different causal effects because they shift treatment for different
margins of the population; the diversity of LATEs is informative
about treatment effect heterogeneity, not a symptom of model
failure.

\begin{example}{Compulsory Schooling Laws vs.\ Distance to College}
\label{w7:ex:late-comparison}
Two classic instruments for years of schooling are (i)~compulsory
schooling laws \citep{angrist1991does} and (ii)~distance to the nearest
college \citep{card1995using}.  The compliers for instrument~(i) are
individuals at the margin of dropping out before the compulsory
leaving age---typically lower-income students.  The compliers
for instrument~(ii) are students deterred from college by
geographic distance---again, disproportionately lower-income.
Both instruments identify the return to schooling for their
respective complier populations, and the LATEs can legitimately
differ even if both instruments are valid.
\end{example}

The dependence of the LATE on the instrument is sometimes described
as a limitation of IV.  It is better understood as a precise
statement about what question is being answered.  A researcher using
compulsory schooling laws is estimating the return to schooling for
students at the compulsory leaving margin; a researcher using
distance to college is estimating the return for students deterred
by geography.  These are different causal questions, and it is
informative---not troubling---that they can yield different answers.

\subsection{The Policy Relevance of LATE}

For many policy questions, the LATE is exactly the right estimand.
If a policy is designed to encourage a subset of the population to
take treatment---for example, an outreach program that reaches
only some potential participants---then the effect on compliers
(those who respond to the encouragement) is precisely what the
policy-maker wants to know.  LATE is most policy-relevant when the
contemplated intervention resembles the instrument, because then
the complier population under the study design is close to the
policy-relevant margin.

When the ATE over the full population is required for policy
analysis, IV alone is insufficient under heterogeneous effects.
Additional assumptions---such as an explicit model of treatment
effect heterogeneity, or a second instrument that identifies
effects for a different subpopulation---are needed to extrapolate
from the LATE to the ATE.


% ============================================================
\section{IV versus Back-Door Adjustment}
\label{w7:sec:compare}
% ============================================================

Having developed both strategies in full, it is instructive to
compare them directly.  The two strategies differ in what they
require, what they identify, and how they fail.

\subsection{A Direct Comparison}

\begin{table}[h]
\centering
\renewcommand{\arraystretch}{1.7}
\begin{tabular}{p{0.18\linewidth} p{0.38\linewidth} p{0.38\linewidth}}
  \toprule
  \textbf{Dimension}
    & \textbf{Back-door / propensity score}
    & \textbf{Instrumental variables} \\
  \midrule
  Core assumption
    & All confounders observed:
      $(Y(0),Y(1)) \indep T \mid X$
    & Valid instrument: relevance, exogeneity,
      exclusion \\[4pt]
  Unobserved confounders
    & Fatal: back-door adjustment fails
    & Permitted: IV routes around~$U$ \\[4pt]
  Estimand
    & ATE, ATT, or ATC depending on design
      and overlap; all coincide under
      homogeneity
    & LATE (compliers only); reduces to a
      common~$\beta$ under homogeneous effects \\[4pt]
  Testability
    & Unconfoundedness is untestable;
      overlap is testable
    & Relevance testable; exogeneity and
      exclusion untestable (in just-identified
      case) \\[4pt]
  Main threat
    & Unmeasured confounder
    & Exclusion restriction violation \\[4pt]
  Identifies ATE?
    & Yes, under strong ignorability
    & Only under homogeneous effects \\[4pt]
  Typical setting
    & Rich administrative or survey data;
      RCT with imperfect compliance
    & Natural experiment; RCT with
      non-compliance; policy change \\
  \bottomrule
\end{tabular}
\renewcommand{\arraystretch}{1}
\end{table}

\subsection{Complementary Failure Modes}

The two strategies have complementary failure modes.

\medskip\noindent\textbf{Back-door adjustment fails} when~$X$ does
not capture all confounders---that is, when an unobserved
variable~$U$ opens a back-door path.  The estimator is then
inconsistent even with infinite data, because the path
$T \leftarrow U \rightarrow Y$ transmits spurious association
that conditioning on~$X$ cannot close.

\medskip\noindent\textbf{IV fails} when the exclusion restriction
is violated---that is, when~$Z$ has a direct effect on~$Y$ beyond
its effect through~$T$.  As derived in
\S\ref{w7:sec:why-matter}, the bias in the Wald estimand is
$\delta/\pi$, amplified by weak instruments.

\medskip\noindent The two failures are in a sense orthogonal:
back-door adjustment requires many observed covariates but tolerates
no unobserved ones, while IV tolerates unobserved confounders but
requires an instrument with no direct effect on the outcome.
Back-door adjustment fails when confounding remains after
conditioning; IV fails when the proposed source of exogenous
variation is not truly exogenous or does not act solely through
treatment.  When designing a study, the choice of identification
strategy should be guided by which assumption is more plausible in
the specific empirical context.

\subsection{When Both Strategies Are Available}

When a valid instrument \emph{and} a sufficient adjustment set~$X$
are both available, comparing the two estimates can be informative,
but disagreement between them does not by itself tell us which
method is wrong.  The two strategies typically target different
estimands---back-door adjustment identifies the ATE or ATT over
the full population or the treated, while IV identifies the LATE
for compliers---and they rely on different identifying assumptions.
Agreement is therefore reassuring under treatment-effect homogeneity
but is neither required nor sufficient otherwise.

The \citet{hausman1978specification} endogeneity test operationalizes
this comparison: under the null hypothesis that~$T$ is exogenous
given~$X$, both the OLS and IV estimators are consistent, and a
large discrepancy is evidence of endogeneity.  Rejection means at
least one strategy is inconsistent; it does not identify which.

\begin{remark}[Interpreting a Hausman Rejection]
A significant Hausman test can arise from two distinct sources:
(a)~back-door adjustment is biased because~$T$ is endogenous, or
(b)~IV is biased because the exclusion restriction is violated or
the instrument is weak.  Even when the test is non-significant,
the two estimates may still differ for the legitimate reason that
they identify different estimands---LATE versus ATE---rather than
because either strategy has failed.  External evidence---instrument
strength diagnostics and institutional arguments for exclusion---is
needed to interpret any discrepancy.
\end{remark}


% ============================================================
\section{Practical Guidance on Defending an IV Design}
\label{w7:sec:practical}
% ============================================================

There is no algorithm for finding a valid instrument; credible
instruments arise from institutional knowledge and careful reasoning
about the data-generating process.  A strong IV design is defended
primarily by institutional knowledge, design logic, and causal
structure; statistical diagnostics are supportive but secondary.

A researcher proposing an instrument should be able to answer
five questions explicitly:

\begin{enumerate}
  \item \textbf{What exactly is the instrument?}  Specify~$Z$
    precisely: its source of variation, the level at which it varies,
    and the population to which it applies.

  \item \textbf{Why does it shift treatment?}  Articulate the causal
    mechanism by which~$Z$ moves~$T$.  Relevance can be verified
    empirically with the first-stage $F$-statistic, but the
    $F$-statistic is a diagnostic for instrument strength, not a
    substitute for a causal account of the $Z \to T$ link.

  \item \textbf{Why is it as-if random relative to latent outcome
    determinants?}  Argue why~$Z$ is unrelated to the unobserved
    causes of~$Y$.  The most credible sources are designed
    randomization (lotteries, randomized encouragement), natural
    experiments with institutional quasi-randomness (policy
    discontinuities, geographic boundaries, biological quirks), and
    shift-share designs \citep{bartik1991benefits,
    goldsmith2020bartik}.  Placebo regressions on pre-determined
    outcomes provide partial---but not definitive---evidence.

  \item \textbf{Why can it affect the outcome only through
    treatment?}  The exclusion restriction is untestable in
    just-identified models, so it must rest on a structural argument
    that no direct path $Z \to Y$ exists.  A useful diagnostic is to
    ask: how large would the direct effect~$\delta$ have to be,
    relative to the first-stage coefficient~$\pi$, to overturn the
    estimated causal effect?  When the first stage is weak, the
    answer is: not very large at all.

  \item \textbf{What population margin does it shift?}  Identify the
    complier population---the units whose treatment status changes
    with~$Z$.  This determines the LATE that is being identified and
    governs the external validity of the estimates for other
    populations or policy margins.
\end{enumerate}


% ============================================================
\section{Applied Example: Charter School Lotteries and the KIPP Lynn Study}
\label{w7:sec:kipp}
% ============================================================

The preceding sections developed the three IV assumptions and their
identification implications in the abstract.  This section grounds
those ideas in a canonical empirical application:
\citet{angrist2012benefits} use the admissions lottery of KIPP Academy
Lynn---a charter middle school in Lynn, Massachusetts---to estimate the
causal effect of KIPP attendance on student achievement.

\subsection{Setting and Instrument}

KIPP (Knowledge Is Power Program) schools follow a ``No Excuses''
model: extended school days, a longer academic year, selective teacher
hiring, and strict behavioral norms.  KIPP Academy Lynn was
substantially oversubscribed beginning in 2005.  Massachusetts law
requires oversubscribed charter schools to select students by lottery,
so the school conducted randomized admissions lotteries from 2005
through 2008.  The lottery offer~$Z_i \in \{0, 1\}$ is the proposed
instrument.  The lottery offer is not the treatment itself but an
exogenous encouragement to enroll, making this a canonical
randomized-encouragement IV design.  The treatment~$s_{igt}$ is the
number of calendar years student~$i$ has spent at KIPP by the time of
test~$(g, t)$---a continuous, endogenous variable, because families
self-select into applying and attending.  The outcome~$Y_{igt}$ is the
student's standardized score on the Massachusetts Comprehensive
Assessment System (MCAS), normalized to mean zero and standard
deviation one within each subject--grade--year cell statewide.

\begin{figure}[h]
\centering
\begin{tikzpicture}[
  node distance = 12mm and 20mm,
  w7obs/.style   = {circle, draw=accent, line width=1pt,
                    minimum size=11mm, font=\small, fill=defbg},
  w7unobs/.style = {circle, draw=backdoorred, line width=1pt, dashed,
                    minimum size=11mm, font=\small, fill=warnbg},
  w7xobs/.style  = {circle, draw=causalgreen, line width=1pt,
                    minimum size=11mm, font=\small, fill=green!8},
  w7arr/.style   = {-{Stealth[length=4.5pt]}, line width=0.9pt,
                    color=accent},
  w7redarr/.style= {-{Stealth[length=4.5pt]}, line width=0.9pt,
                    color=backdoorred, dashed},
  w7grnarr/.style= {-{Stealth[length=4.5pt]}, line width=0.9pt,
                    color=causalgreen},
  w7lbl/.style   = {font=\footnotesize\color{isublue}}
]

%% Nodes
\node[w7obs]                              (Z) {$Z$};
\node[w7obs, right=of Z]                  (S) {$S$};
\node[w7obs, right=of S]                  (Y) {$Y$};
\node[w7unobs, above=of S, xshift=10mm]   (U) {$U$};
\node[w7xobs,  below=of S, xshift=10mm]   (X) {$X$};

%% Causal edges
\draw[w7arr]    (Z) -- node[above, w7lbl]{\scriptsize relevance} (S);
\draw[w7arr]    (S) -- node[above, w7lbl]{\scriptsize $\rho$}    (Y);
\draw[w7redarr] (U) -- (S);
\draw[w7redarr] (U) -- (Y);
\draw[w7grnarr] (X) -- (S);
\draw[w7grnarr] (X) -- (Y);

%% Absent-arrow annotations
\node[font=\scriptsize\color{backdoorred}, above=3mm of Z,
      xshift=2mm, align=center]
  {$Z \not\leftarrow U$\\[-1pt](exogeneity)};
\draw[backdoorred!40, densely dashed, very thin,
      -{Stealth[length=3pt]}, bend left=40]
  (Y) to node[below, font=\scriptsize\color{backdoorred!70}]
  {no $Z \to Y$\;(exclusion)} (Z);

\end{tikzpicture}

\medskip
{\small \textbf{Variable key.}
$Z$ = lottery offer (instrument);
$S$ = years enrolled at KIPP (endogenous treatment);
$Y$ = MCAS test score (outcome);
$U$ = unobserved determinants of both enrollment and achievement
  (e.g.\ parental engagement, student motivation);
$X$ = observed baseline covariates (race, LEP/SPED status, prior
  test scores).
Blue arrows denote causal effects.  Dashed red arrows are paths
through the unobserved confounder.  Green arrows are observed
covariate effects.  Two absent arrows encode the IV assumptions:
there is no $U \to Z$ or $Z \leftarrow U$ path (exogeneity), and
there is no direct $Z \to Y$ arrow (exclusion).}
\caption{Causal DAG for the KIPP Lynn lottery IV design.
The lottery offer~$Z$ shifts enrollment~$S$ (relevance), is
unrelated to the unobserved confounder~$U$ by random assignment
(exogeneity), and has no direct path to the test score~$Y$
(exclusion).}
\label{w7:fig:kipp-dag}
\end{figure}

\subsection{Mapping the Three Assumptions to the KIPP Context}

Figure~\ref{w7:fig:kipp-dag} displays the causal structure.  Each
IV assumption corresponds to a specific feature of the DAG; the
three paragraphs below trace that correspondence.

\textbf{Relevance.}  The lottery offer must shift years of KIPP
attendance.  Lottery winners were offered a seat and about 80~percent
accepted; losers rarely enrolled elsewhere at KIPP.  The first-stage
regression of~$s_{igt}$ on~$Z_i$ (with year and grade controls) yields
a coefficient of approximately~$1.2$: at the time of each MCAS exam,
lottery winners had spent about 1.2~more years at KIPP than lottery
losers.  The first-stage $F$-statistic is far above conventional
thresholds.  The first stage is less than the theoretical maximum
because compliance is partial (some winners do not attend; some
losers find entry through later cohorts or attrition slots) and the
follow-up window captures students at different points in their KIPP
tenure.

\textbf{Exogeneity.}  The lottery offer must be independent of all
potential outcomes.  Because offer status was determined by randomly
drawn lottery-sequence numbers, this condition holds by design within
each application cohort---an especially strong basis for exogeneity
compared to most observational IV applications.  The authors verify it
empirically: a joint test of covariate balance across lottery winners
and losers yields a $p$-value of~$0.615$ for demographic
characteristics and baseline test scores, consistent with the null of
no pre-lottery differences (Table~1 of \citealt{angrist2012benefits}).
Pre-lottery variables are used for the balance check precisely because
post-lottery variables such as LEP or SPED classification may
themselves be affected by school attended.

\textbf{Exclusion.}  The lottery offer must affect test scores only
through KIPP attendance, not through any direct channel.  Because the
offer merely provides access to a school---it does not itself deliver
instruction---this restriction is institutionally plausible.  One
potential violation is a discouragement effect: losing the lottery
might demoralize students, suppressing their achievement regardless of
where they enroll.  The authors address this by noting that the scores
of lottery losers are typical of demographically comparable students in
Lynn, which is inconsistent with large discouragement effects.

\subsection{First Stage, Reduced Form, and the 2SLS Estimand}

The equations below are used here only to illustrate the components
of the Wald/2SLS logic---first stage, reduced form, and their ratio.
Formal estimation theory for 2SLS, including asymptotic variance and
cluster-robust inference, is deferred to Chapter~12.

The structural equation for test scores is
\begin{equation}
  y_{igt} = \alpha_t + \beta_g + \sum_j \delta_j d_{ij}
            + \gamma' X_i + \rho\, s_{igt} + \varepsilon_{igt},
  \label{w7:eq:kipp-structural}
\end{equation}
where $\alpha_t$ and $\beta_g$ are year-of-test and grade-of-test
fixed effects, $d_{ij}$ are application-cohort dummies
(cohort membership determines the probability of winning, so cohort
is an essential control), $X_i$ is a vector of baseline demographics,
and $\rho$ is the causal effect of interest per year at KIPP.  The
first-stage equation is
\begin{equation}
  s_{igt} = \lambda_t + \kappa_g + \sum_j \mu_j d_{ij}
            + \Gamma' X_i + \pi Z_i + \eta_{igt}.
  \label{w7:eq:kipp-firststage}
\end{equation}
The model is just-identified (one excluded instrument per endogenous
variable), so the 2SLS estimator of~$\rho$ equals the ratio of the
reduced-form coefficient on~$Z_i$ to the first-stage coefficient~$\pi$.

Estimates from \citet{angrist2012benefits} (their Table~4, the
specification with demographic and baseline-score controls) are
summarized in Table~\ref{w7:tab:kipp}.

\begin{table}[ht]
\centering
\caption{IV estimates of KIPP Lynn attendance on MCAS scores
         (per year at KIPP).
         Source: \citet{angrist2012benefits}, Table~4.}
\label{w7:tab:kipp}
\begin{tabular}{lccc}
\hline\hline
Subject & First stage & Reduced form & 2SLS \\
\hline
Math & $1.221$ & $0.430$ & $0.352$ \\
     & $(0.068)$ & $(0.067)$ & $(0.053)$ \\
ELA  & $1.228$ & $0.164$ & $0.133$ \\
     & $(0.068)$ & $(0.073)$ & $(0.059)$ \\
\hline\hline
\multicolumn{4}{p{0.82\linewidth}}{\small Standard errors clustered at
the student level in parentheses.  $N = 833$ student-by-test
observations for both subjects.  All regressions include year-of-test,
grade-of-test, and application-cohort fixed effects, plus demographic
and baseline-score controls.}
\end{tabular}
\end{table}

Each year at KIPP raises math scores by approximately $0.35$ standard
deviations and ELA scores by approximately $0.13$ standard deviations.
The reduced-form estimate for math ($0.43\sigma$) is larger than the
2SLS estimate ($0.35\sigma$) because the first stage exceeds~$1$:
lottery winners accumulated somewhat more than one additional year at
KIPP per unit of follow-up time.  The reduced form regression, the
2SLS estimator, and its asymptotic theory are all developed in
Chapter~12.

\begin{tcolorbox}[remarkstyle, title={Remark: Treatment Heterogeneity and the Weighted-Average Interpretation}]
When the treatment is continuous, \citet{angrist1995two} show that
the IV estimand is a weighted average of the marginal causal effects
of each additional unit of treatment, with weights proportional to
the first-stage effect of the instrument at each dose level.  In the
KIPP setting, $\rho$ is therefore a weighted average of the
per-year effect across the different years a student might spend at
KIPP, rather than the effect of a single fixed dose.  The authors
note this interpretation explicitly and adopt the conservative coding
convention that partial years count as full years, which slightly
attenuates the 2SLS estimates.
\end{tcolorbox}

\subsection{LATE Interpretation}

In the KIPP lottery, compliance is partial: some lottery winners do
not enroll and some lottery losers eventually find entry.  The
relevant compliance types are therefore defined relative to offer
status.  \emph{Compliers} are students who attend KIPP if and only
if they win the lottery.  Always-takers and never-takers are ruled
out at the extensive margin by the lottery's control over seat
access, though they reappear at the intensive margin (years of
enrollment).  The 2SLS estimand~$\rho$ is therefore the LATE: the
average per-year treatment effect for the \emph{lottery compliers},
not for all applicants and not for all students in Lynn.

An important feature of this LATE is that it may differ from the
effect the school would have on a randomly selected Lynn student who
was not part of the applicant pool.  KIPP applicants already had
parents motivated enough to apply.  The authors find that KIPP
applicants have baseline test scores slightly \emph{lower} than the
LPS average (Table~1 of \citealt{angrist2012benefits}), so the
applicant pool is not positively selected on prior achievement, but
it may still be selected on parental engagement in ways that affect
how students respond to the KIPP environment.

\subsection{Treatment Effect Heterogeneity}

Subgroup analysis reveals that the LATE varies substantially across
observable student characteristics.  Reading gains ($\approx 0.12\sigma$
overall) are driven almost entirely by students classified as having
limited English proficiency (LEP, $\approx 0.43\sigma$) and special
education needs (SPED, $\approx 0.27\sigma$); non-LEP, non-SPED
students show negligible ELA gains.  Math effects are large and
positive across all subgroups but are largest for LEP and
lower-achieving students.  An interaction model that adds the product
of baseline score with years at KIPP---identified by including
$Z_i \times \text{baseline score}$ as a second instrument---yields
a significantly negative interaction term in both subjects: each
additional standard deviation of baseline disadvantage is associated
with an additional $0.08$--$0.17\sigma$ gain per year at KIPP.

These findings connect directly to Section~\ref{w7:subsec:diff-instr}:
were we to define separate subgroup-specific instruments (LEP-lottery
and non-LEP-lottery), each would identify a distinct subpopulation
LATE.  The overall 2SLS estimate is a weighted average of these
subgroup LATEs, with weights proportional to each subgroup's share of
the complier population.  This heterogeneity is not a nuisance
detail; it is exactly why IV estimates must be interpreted together
with the margin of compliance they capture.

\begin{tcolorbox}[remarkstyle, title={Remark: Lottery Design and Assumption Credibility}]
The KIPP study illustrates a general principle: instruments derived
from \emph{designed} randomization---lotteries, random assignment,
randomized encouragement---provide the most transparent basis for
the exogeneity assumption, because the distribution of~$Z$ is known
by construction.  The exclusion restriction still requires
institutional argument (here, that the lottery offer affects only
schooling, not morale or parental behavior in other dimensions), but
the exogeneity assumption is not merely plausible---it is guaranteed
by the randomization protocol.  This is why lottery-based IV designs
occupy a privileged position in the program evaluation literature
\citep{angrist2009mostly}.
\end{tcolorbox}

% ============================================================
\section{Summary}
\label{w7:sec:summary}
% ============================================================

\begin{enumerate}
  \item \textbf{IV identifies effects from exogenous treatment
    variation.}  When back-door adjustment fails because an
    unobserved~$U$ creates a path $T \leftarrow U \rightarrow Y$,
    a valid instrument~$Z$ identifies the causal effect by exploiting
    only the component of treatment variation that~$Z$ induces.  IV
    does not block the confounding path---it avoids it by isolating
    exogenous variation and ignoring the rest.

  \item \textbf{Three assumptions, ordered by testability.}
    Relevance---that $Z$ causally shifts $T$---can be assessed with
    the first-stage $F$-statistic, though the $F$-statistic is a
    sample diagnostic, not the assumption itself.  Exogeneity and
    exclusion are primarily substantive: exogeneity requires that~$Z$
    carries no information about latent outcome determinants other
    than through~$T$; exclusion requires that~$Z$ has no direct
    effect on~$Y$.  Both must be defended by institutional knowledge,
    design logic, and causal structure---statistical diagnostics are
    supportive but secondary.  Each violated assumption produces a
    distinct, quantifiable bias in the Wald estimand, amplified by
    weak instruments (\S\ref{w7:sec:why-matter}).

  \item \textbf{Framework~1: homogeneous-effect SEM
    $\Rightarrow$ Wald identifies $\beta$.}
    Under constant treatment effects and the linear structural
    model, the Wald estimand identifies the single causal
    parameter~$\beta$, which equals the ATE, ATT, and LATE
    simultaneously.  The identification follows from a three-step
    do-calculus argument using the reduced form and first stage
    (\S\ref{w7:sec:linear}).

  \item \textbf{Framework~2: heterogeneity + monotonicity
    $\Rightarrow$ Wald identifies LATE.}
    Under heterogeneous treatment effects and the monotonicity
    assumption, the Wald estimand identifies the average treatment
    effect for \emph{compliers} only---those whose treatment status
    changes with the instrument, a latent subgroup defined by
    potential treatment statuses $(T(0), T(1))$, not by observable
    characteristics (\S\ref{w7:sec:late}--\ref{w7:sec:interpret}).

  \item \textbf{Different instruments identify different effects.}
    The LATE depends on the instrument through the complier
    population it selects.  Different instruments can legitimately
    identify different causal effects because they shift treatment
    for different margins of the population.  This is informative
    about treatment effect heterogeneity, not a contradiction.
    LATE equals ATE only under additional structure, most notably
    treatment-effect homogeneity (\S\ref{w7:sec:interpret}).

  \item \textbf{IV versus back-door adjustment.}  The two
    strategies have complementary failure modes: back-door adjustment
    fails when confounders are unobserved; IV fails when the exclusion
    restriction is violated or the instrument is not truly exogenous.
    IV identifies the LATE (not the ATE) under heterogeneous effects;
    back-door adjustment identifies the ATE or ATT.  They are tools
    for different identification problems, not competitors
    (\S\ref{w7:sec:compare}).

  \item \textbf{Applied example: the KIPP Lynn lottery.}
    \citet{angrist2012benefits} use a randomized admissions lottery
    as a canonical randomized-encouragement instrument for years of
    charter school attendance.  The lottery guarantees exogeneity by
    design; exclusion is defended institutionally; relevance is
    confirmed by a first-stage coefficient of~$1.2$.  The estimates
    identify the LATE for lottery compliers---$0.35\sigma$ per year
    in math, $0.13\sigma$ in ELA---with the largest gains for LEP,
    SPED, and low-baseline-score students, illustrating that the
    margin of compliance shapes which LATE is recovered
    (\S\ref{w7:sec:kipp}).

  \item \textbf{Estimation deferred to Chapter~12.}  This chapter
    establishes what IV identifies and under what assumptions.
    How the Wald ratio is estimated from finite data---the reduced
    form regression, two-stage least squares, asymptotic inference,
    and overidentification tests---is the subject of Chapter~12.
\end{enumerate}

% ============================================================
\section*{Problems}
\addcontentsline{toc}{section}{Problems}
% ============================================================

\begin{enumerate}

\item \textbf{The three IV assumptions in three languages.}
  Consider the DAG:
  $\{Z \to T,\; T \to Y,\; U \to T,\; U \to Y,\;
   X \to T,\; X \to Y,\; X \to Z\}$ with $U$ unobserved.
  \begin{enumerate}
    \item List all back-door paths from $T$ to $Y$.  Does $X$ alone
      satisfy the back-door criterion?  Explain.
    \item Verify the three IV assumptions using $d$-separation:
      (i)~Relevance: show $Z$ and $T$ are not $d$-separated in
      $\Gcal$.
      (ii)~Exogeneity: show $Z \indep U \mid X$ in $\Gcal$.
      (iii)~Exclusion: show $Y \indep Z \mid T, X$ in
      $\Gcal_{\overline{T}}$.
    \item Now add the arrow $Z \to Y$ to the DAG.  Which IV
      assumption is violated?  Show explicitly which step of the
      Wald derivation in \S\ref{w7:sec:linear} breaks down.
    \item Translate each of the three IV assumptions into the
      structural language: write the equations for $T$ and $Y$
      and identify which coefficient restriction corresponds to
      each assumption.
  \end{enumerate}

\item \textbf{Bias under assumption violations.}
  Let $Y = \beta T + \varepsilon$ and $T = \pi Z + \eta$ with
  $\E[\varepsilon \mid Z] = 0$ and $\pi \neq 0$.
  \begin{enumerate}
    \item Starting from $\E[Y \mid Z{=}1] - \E[Y \mid Z{=}0]$,
      substitute the structural equation for $Y$ and simplify.
      What role does exogeneity play?
    \item Show that $\E[T \mid Z{=}1] - \E[T \mid Z{=}0] = \pi$
      in the linear first-stage model.  What role does relevance
      play?
    \item Derive the Wald estimand and confirm it equals $\beta$.
    \item Now suppose the exclusion restriction fails and
      $Y = \beta T + \delta Z + \varepsilon$ with $\delta \neq 0$.
      Derive the probability limit of the Wald estimator.  Confirm
      the bias formula from \S\ref{w7:sec:why-matter}.
    \item Suppose instead that exogeneity fails:
      $\E[\varepsilon \mid Z] = cZ$ for some constant $c \neq 0$.
      Derive the probability limit of the Wald estimator and
      express the bias in terms of~$c$ and~$\pi$.  Compare
      the structure of this bias with the exclusion violation bias.
  \end{enumerate}

\item \textbf{Order, rank, and the limits of overidentification.}
  Consider a model with one endogenous variable~$T$ and two
  instruments $Z_1$ and $Z_2$, both satisfying exogeneity and
  exclusion.
  \begin{enumerate}
    \item State the order condition and verify it is satisfied.
    \item State the rank condition.  What would it mean
      geometrically if the rank condition failed---i.e., if $Z_1$
      and $Z_2$ were perfectly collinear in the first-stage
      regression?
    \item Explain intuitively why having two valid instruments
      rather than one should improve estimation precision.
    \item Now suppose $Z_1$ is valid but $Z_2$ violates the
      exclusion restriction.  The Sargan--Hansen $J$-test is
      applied.  Under what conditions does the $J$-test have
      power to detect $Z_2$'s invalidity?  Under what conditions
      does the test fail to detect it?
    \item Why does passing the $J$-test not confirm that both
      $Z_1$ and $Z_2$ are valid?  Give a concrete example of a
      situation in which both instruments are invalid and the
      $J$-test has no power.
  \end{enumerate}

\item \textbf{Compliance types and the LATE.}
  In a binary instrument, binary treatment study, suppose the
  population has the following composition: 30\% compliers with
  average treatment effect $\tau_c = 6$; 25\% always-takers with
  $\tau_a = 3$; 45\% never-takers with $\tau_n = 1$; no defiers.
  \begin{enumerate}
    \item Compute $P(\text{complier}) = \E[T \mid Z{=}1] -
      \E[T \mid Z{=}0]$.
    \item Compute the ATE as a weighted average of $\tau_c$,
      $\tau_a$, $\tau_n$ with appropriate weights.
    \item The Wald estimand equals $\tau_c = 6$.  By how much
      does this overstate the ATE, and why?
    \item A second study uses a different binary instrument $Z'$
      with a complier population of 50\% and a LATE of~2.  Is
      this contradictory?  What can you infer about the relative
      treatment effect in the two complier populations?
    \item Explain, using compliance type language, why the
      denominator of the Wald estimand equals
      $P(\text{complier})$.
  \end{enumerate}

\item \textbf{The exclusion restriction: plausibility and
  violations.}
  Evaluate the exclusion restriction for each of the following
  proposed instruments.  For each, (i)~state whether the
  restriction is plausible and why; (ii)~describe a specific
  mechanism by which it could be violated; and (iii)~assess
  whether the violation would bias the IV estimate upward or
  downward.
  \begin{enumerate}
    \item \textbf{Instrument}: rainfall in the home region of a
      politician, used as an instrument for government
      infrastructure spending.  \textbf{Outcome}: local economic
      growth.
    \item \textbf{Instrument}: distance to the nearest hospital,
      used as an instrument for hospital admission.
      \textbf{Outcome}: 30-day mortality.
    \item \textbf{Instrument}: a randomly assigned financial
      incentive to enroll in a health screening program, used as
      an instrument for screening uptake.  \textbf{Outcome}:
      health status two years later.
    \item \textbf{Instrument}: lottery number in the Vietnam-era
      draft lottery, used as an instrument for military service.
      \textbf{Outcome}: lifetime earnings.
      [\emph{This is the Angrist~(1990) study; discuss why this
      instrument is widely regarded as satisfying the exclusion
      restriction.}]
  \end{enumerate}

\item \textbf{IV versus back-door adjustment.}
  A researcher studies the effect of job training~($T$) on
  earnings~($Y$).  Two strategies are available: (A)~a rich set
  of pre-treatment covariates~$X$ (education, age, prior earnings,
  region) and a propensity-score estimator; (B)~a lottery that
  randomly selected units to be offered training (not required to
  attend), used as an instrument~$Z$.
  \begin{enumerate}
    \item Under what assumption does strategy~(A) identify the
      ATE?  What specific unobserved variable would most plausibly
      violate this assumption?
    \item Strategy~(B) identifies a LATE.  Describe the complier
      population in words.  Is the LATE likely to be larger or
      smaller than the ATE in this setting?  Explain.
    \item Both strategies are implemented and yield estimates
      of \$1{,}800 and \$2{,}400 per year, respectively.
      Describe a Hausman-type test that uses both estimates.
      Under what null hypothesis does the test have an
      approximate $\chi^2$ distribution?
    \item If the two estimates differ significantly, which
      strategy would you trust more and why?  What additional
      evidence would help distinguish the two explanations
      (endogeneity bias in~(A) versus LATE $\neq$ ATE in~(B))?
  \end{enumerate}

\end{enumerate}

% Bibliography (standalone only)
\ifSubfilesClassLoaded{%
  \bibliographystyle{plainnat}
  \bibliography{causal_inference}
}{}

\end{document}